\documentclass[10pt,conference]{IEEEtran}

\IEEEoverridecommandlockouts
% The preceding line is only needed to identify funding in the first footnote. If that is unneeded, please comment it out.
\usepackage{cite}
\usepackage{amsmath,amssymb,amsfonts}
\usepackage{algorithmic}
\usepackage{graphicx}
\usepackage{textcomp}
\usepackage{booktabs}
\usepackage{multirow}
\usepackage{xcolor}
\usepackage{url}
\usepackage{listings}
\usepackage{tcolorbox}
\usepackage{balance}
\newtcolorbox{rqanswer}[1]{colback=gray!7, colframe=gray!55, boxrule=0.4pt,
  arc=1pt, left=4pt, right=4pt, top=3pt, bottom=3pt, before skip=6pt,
  after skip=6pt, fontupper=\small, before upper={\textbf{Answer to #1.}~}}
\def\BibTeX{{\rm B\kern-.05em{\sc i\kern-.025em b}\kern-.08em
    T\kern-.1667em\lower.7ex\hbox{E}\kern-.125emX}}

\newcommand{\red}[1]{\textcolor{red}{#1}}
\newcommand{\orange}[1]{\textcolor{orange}{#1}}


%\raggedbottom

\begin{document}

%\title{Predicting Code Correctness under Distribution Shift: A Study of Internal and Causal LLM Signals}

\title{Which Signals of Code Correctness Survive Distribution Shift? Hidden-State and Causal Predictors of Program and Verdict Correctness}

% Can the Judge Be Trusted? Execution-Free Correctness Prediction for LLM-Generated Code under Distribution Shift
% Where Execution-Free Correctness Predictors Break: Benchmark and Language Shift in LLM Code Generation
% Causal Subspaces vs. Probes: Predicting the Correctness of LLM-Generated Code under Distribution Shift
% Which Signals Tell When an LLM's Code Judgment Fails? Comparing Correctness Predictors under Language and Benchmark Shift
% Causal Subspaces Transfer Across Languages, Not Across Tasks: Comparing Correctness Predictors for LLM-Generated Code
% Predicting Code Correctness under Distribution Shift: A Study of Internal and Causal Signals
% Out-of-Distribution Prediction of Code Correctness and LLM Verdict Reliability
% When Can We Trust an LLM's Code Judgments? Execution-Free Prediction under Distribution Shift
% When Can We Trust an LLM's Verdict on Its Own Code? Execution-Free Prediction under Benchmark and Language Shift
% Which Signals of Code Correctness Survive Distribution Shift? Confidence, Probes, and Causal Subspaces for LLM-Generated Code

%\thanks{Identify applicable funding agency here. If none, delete this.}


\author{\IEEEauthorblockN{%Anonymous Author(s) %
Viola Campos, Thomas Klassert, and Adrian Ulges
}
\IEEEauthorblockA{\textit{RheinMain University of Applied Sciences}\\
Wiesbaden, Germany \\
\{viola.campos, adrian.ulges\}@hs-rm.de} thomas.klassert@student.hs-rm.de
 % \and
 % \IEEEauthorblockN{Thomas Klassert}
 % \IEEEauthorblockA{\textit{dept. name of organization (of Aff.)} \\
 % \textit{RheinMain University of Applied Sciences}\\
 % Wiesbaden, Germany \\
 % email address or ORCID}
 % \and
 % \IEEEauthorblockN{Adrian Ulges}
 % \IEEEauthorblockA{\textit{dept. DCSM} \\
 % \textit{RheinMain University of Applied Sciences}\\
 % Wiesbaden, Germany \\
 % email address or ORCID}
}


\maketitle

\begin{abstract}
% Matched to the introduction (2026-10-29); previous version below as comment.
Large language models increasingly judge code, %including their own, 
yet without external feedback such judgments are unreliable, and verifying them by
executing tests or by extended reasoning is costly. %Execution-free correctness predictors offer a cheap signal from the model itself, from its output probabilities, the agreement among sampled solutions, or its hidden states, that can decide when a judgment needs costlier verification. 
Instead, execution-free correctness predictors use signals from the model's output probabilities, sampled solutions, or hidden states to identify judgments that warrant further verification.
However, 
these predictors 
%trained on hidden states 
may degrade under distribution
shifts 
%an issue that remains underexplored for code, 
when generalizing across projects and languages. 

We compare six execution-free predictors from four families on Qwen3-Coder-30B and gpt-oss-20b under benchmark shift between LiveCodeBench and BigCodeBench and language shift across Python, C++ and Java.
We evaluate two targets: whether the model's own verdict on a program can be trusted, and whether the program itself is correct.
Besides established approaches, we adapt two novel causal localization methods to code, which have been shown to be more robust out-of-distribution on natural language tasks. 

%For verdict correctness on Qwen3-Coder, causal localization probes improve robustness by 0.11 in distribution and by 0.09 under language shift, significantly outperforming the best non-causal predictor by 0.04 AUROC on average.
For verdict correctness, causal localization improves robustness to language shift. For Qwen3-Coder a probe on the causally localized subspace stays 0.09 AUROC ahead of output confidence under language shift (0.11 in distribution), significantly outperforming the best non-causal predictor by 0.04 AUROC on average.
%For verdict correctness, causal localization probes improves robustness to language shift. For Qwen3-Coder a probe on the localized subspace retains a 0.09 AUROC advantage over output confidence, compared with 0.11 in distribution, while a full-hidden-state probe retains only 0.03. The localized-subspace probe also significantly outperforms the best non-causal predictor by 0.04 AUROC on average. 
This advantage, however, does not generalize consistently to gpt-oss or to benchmark shift. For program correctness, predictors transfer more successfully between benchmarks, largely by identifying problems on which the model tends to fail. However, their ability to rank candidates within the same problem is limited. % Agreement across sampled solutions offers no consistent advantage over a single-pass signal.

\end{abstract}

\begin{IEEEkeywords}
code correctness prediction, uncertainty estimation, LLM-as-a-judge, distribution shift, causal interpretability
\end{IEEEkeywords}


\section{Introduction}
\label{sec:intro}
% Draft 2026-10-28, shortened from 2.4 to about 1.2 columns (long version
% below as comments). Framing: (1) causal localization pays off for the verdict
% target, (2) general assessment of generalization (Adrian's comment on the
% abstract). Motivation: cheap signals gate costly verification (literature
% checked 2026-10-28). Numbers: RQ answer boxes in experiments.tex,
% paired_shift.csv (fold=mean).
LLMs are increasingly used to judge code, including their own. They determine whether a candidate solves
its task, whether a change is safe to accept, and which of several candidates to
keep. Without external feedback, such judgments are unreliable. Models rarely
improve by verifying their own outputs~\cite{Huang2024SelfCorrect,
Kamoi2024SelfCorrection}, favor their own generations~\cite{Panickssery2024Self},
and remain inconsistent as code judges~\cite{Jiang2025CodeJudgeBench}. Reliable
verification is costly: executing tests requires an environment and a test
suite~\cite{Raghavendra2026Rubrics}, and reasoning about a candidate involves generating many additional
tokens. A cheap signal indicating when a judgment can be trusted could gate
these costlier checks rather than replace them~\cite{Rakasi2026Uncertainty,
jung2025trust}. Execution-free correctness predictors provide such a signal from
the model itself, such as the probability of its
answer~\cite{Spiess2025Calib, Campos2026Multicalib}, the agreement among sampled
solutions~\cite{Chen24Eigenscore, shi2026codemorethantext,
bouchard2026functional}, or its hidden states~\cite{bui2025openia,
Vu2026Autoprobe, ribeiro2026repe}. Probes on hidden states are particularly
cheap, since they reuse activations the model computes anyway, and are already
deployed as first-stage monitors~\cite{McKenzie2025HighStakes,
Kramar2026Probes}. Such a predictor serves two purposes: it can estimate whether
a program is correct, and whether the model's own verdict on it can be trusted.

Although predictors trained on hidden states are accurate in distribution, their
acknowledged weakness is distribution shift~\cite{Kramar2026Probes,
OrgadTGRSKB25}. For code this is critical, as a predictor fitted on one benchmark
or programming language will usually be applied to other projects and languages. For natural-language question answering, Huang et
al.~\cite{Huang2025Internal} proposed causal localization as a remedy: they use
counterfactual interventions, which copy part of the hidden state from one
input into another and observe how the model's answer changes, to find the
subspace that robustly determines the answer %, and report that predictors built on it
%are more robust 
under shift. Whether this promise holds for code correctness assessment 
%, and how the established predictors compare under shift, 
remains open. Instead, prior work on code has considered only a subset of predictor families
and shift settings, and has largely focused on predicting program
correctness~\cite{ribeiro2026robustness, he2026codecircuit, huang2025pttrust,
gros2025localized}. To our knowledge, none includes causal methods or asks whether the model's
own verdict can be trusted.

We therefore ask \emph{which execution-free signals of code correctness
survive a change of benchmark or programming language, and what explains those
that do}. We adapt two causal methods of Huang et al., \emph{counterfactual
simulation} and \emph{value probing}, to code and compare them with output
confidence, a probe on the hidden state, LAT~\cite{ribeiro2026repe} and
EigenScore~\cite{Chen24Eigenscore}. To do so, we use Qwen3-Coder-30B and gpt-oss-20b for candidate generation and as judges; we apply shifts between LiveCodeBench and BigCodeBench
and between Python, C++ and Java. We then answer four research questions: \textbf{RQ1} asks how well the methods tell
whether the model's verdict can be trusted, the target the causal methods were
designed for; \textbf{RQ2} asks which part of the causal methods carries their
signal; \textbf{RQ3} asks how well the methods predict whether the program is
correct, the target of most prior work; and \textbf{RQ4} asks whether
agreement across samples, which requires no supervised training, tells better than a single
pass which problems the model will solve. 

% Findings without numbers (2026-10-29): the numbers are in the abstract and
% RQ1/RQ3. Previous numeric version: 0.09 of 0.11 kept vs 0.03 (verdict,
% Qwen); +0.04 over the best non-causal method; program 0.05 of 0.12 (Qwen),
% 0.02 of 0.16 (gpt-oss).
Our results show that causal localization pays off where it was designed to:
under language shift, a probe on the localized subspace is the most robust
predictor of Qwen3-Coder's verdict, significantly ahead of every non-causal
method, and the learned subspace itself, not the intervention, carries this
gain. On gpt-oss, no method stands out. Beyond the causal methods,
generalization depends on the target and the type of shift: no verdict
predictor transfers reliably between benchmarks, program correctness
transfers largely because the predictors recognize difficult problems, and
agreement across samples gives no consistent advantage over a single pass.

% Under language shift, the methods differ in how much of their lead over the
% untrained output confidence they keep. On the verdict target, causal
% localization keeps most of it: for Qwen3-Coder, a probe on the localized
% subspace stays 0.09 AUROC ahead of confidence (0.11 in distribution), while a
% probe on the full hidden state keeps only 0.03. The subspace probe is significantly ahead of
% the best non-causal method ($+0.04$ AUROC on average) and ahead in each of the
% three data partitions. The learned subspace itself carries this gain, so a
% simple probe on it suffices; random subspaces of the same size do not provide
% it. On gpt-oss, the causal methods are level with the best non-causal method.
% On program correctness, the trained methods keep at most half of their lead
% (0.05 of 0.12 for Qwen3-Coder, 0.02 of 0.16 for gpt-oss). Between benchmarks,
% the picture reverses: no method transfers verdict prediction reliably, whereas
% program correctness transfers, largely because the methods recognize which
% problems the model fails on; within one problem every method is weak.
% Agreement across samples gives no consistent advantage over a single pass.


Our contributions are:
\begin{itemize}
  \item To our knowledge, the first comparison of execution-free correctness
  predictors for code across four families, including causal methods, under
  both benchmark and language shift and for two targets.
  \item An adaptation of causal localization~\cite{Huang2025Internal} to code
  correctness prediction.
  \item An analysis showing that the learned subspace carries the advantage of
  the causal methods, and a distinction between candidate-level prediction and
  problem-level reliability: pooled AUROC must be read alongside a predictor's
  ability to rank the candidates of one problem.
\end{itemize}

% ---------------------------------------------------------------------------
% Long version (2.4 columns, with Viola's edits; replaced 2026-10-28). Parts
% cut and where they live now: self-preference \cite{Panickssery2024Self}
% (dropped); "two purposes" (RQ1/RQ3); RQ display blocks (inline); probe loses
% three times as much (RQ1); protocol sentence (Study Design); DAS training
% measure does not indicate usefulness (RQ2); problem-only probe (RQ3);
% contribution-2 detail (Methods).
%
% LLMs increasingly judge code, including their own They determine whether a
% candidate solves its task, whether a change is safe to accept, and which of
% several candidates to keep. Without external feedback, such judgments are
% unreliable. Models rarely improve by verifying their own
% outputs~\cite{Huang2024SelfCorrect, Kamoi2024SelfCorrection}, favor their own
% generations~\cite{Panickssery2024Self}, and remain inconsistent as code
% judges~\cite{Jiang2025CodeJudgeBench}. Reliable verification is costly:
% executing tests requires an environment and a test
% suite~\cite{Raghavendra2026Rubrics}, and reasoning about a candidate involves
% generating many additional tokens. A cheap signal indicating when a judgment
% can be trusted could gate these costlier checks rather than replace
% them~\cite{Rakasi2026Uncertainty, jung2025trust}. Execution-free correctness
% predictors provide such a signal from the model itself, ... Such a predictor
% serves two purposes: it can estimate whether a program is correct, and
% whether the model's own verdict on it can be trusted.
%
% Although predictors trained on hidden states are accurate in distribution,
% their acknowledged weakness is distribution shift. For code this is
% critical: a predictor fitted on one benchmark or programming language will
% be applied to other projects and languages. For natural-language question
% answering, Huang et al. proposed causal localization as a remedy. They use
% counterfactual interventions to identify a subspace of the hidden state that
% causally determines the model's answer. According to the authors, a
% predictor based on this subspace is more robust under distribution shift.
% Whether this holds for code, and how the established predictors compare
% under shift, is open. Prior studies of execution-free correctness prediction
% for code examine one or two families of predictors, under one type of shift
% and for program correctness only. To our knowledge, no study has compared
% execution-free predictors across method families under both a change of
% benchmark and programming language. None includes causal methods or asks
% whether the model's own verdict can be trusted.
%
% We therefore ask: which execution-free signals of code correctness survive a
% change of benchmark or programming language, and what explains those that
% do? ... Four research questions follow from the overarching one. The causal
% methods model the judge's decision, so we begin with the target they were
% designed for.
% RQ1: How well can execution-free methods tell whether the model's own
%   verdict can be trusted, in distribution and under shift?
% RQ2: Which part of the causal methods carries their signal? Where they are
%   helpful, we separate the learned subspace from the intervention and
%   compare it with random subspaces.
% RQ3: How well do the methods predict whether the program is correct? This is
%   the target of most prior work related to code.
% RQ4: Does agreement across samples indicate better than a single pass which
%   problems the model will solve? Agreement requires no training and
%   therefore cannot degrade under shift.
%
% Findings: ... In contrast, a probe on the full hidden state loses three
% times as much. The advantage lies in the learned subspace itself: a simple
% probe on it is as robust as the full causal method, and thus as cheap as any
% probe, while random subspaces of the same size do not provide this
% improvement. ...
%
% Contributions: (1) To our knowledge, the first broad comparison ... No
% reported number depends on a layer chosen on the same data, the language
% experiments are repeated on three partitions of the problems, and methods are
% compared with paired significance tests. (2) Novel adaptation of causal
% localization to code correctness prediction. One subspace per layer is
% trained on candidates the model judges correctly, serving both targets. (3)
% Analysis: the learned subspace, not the intervention, carries the gain, so a
% simple probe on it suffices; tracing how well the subspace is learned does
% not indicate whether it is useful. (4) Two findings for the program
% correctness target: pooled AUROC largely reflects which problems the model
% solves, since a probe that has seen only the problem matches one that has
% seen the candidate; and under language shift, trained methods lose most of
% their lead over output confidence.

% ---------------------------------------------------------------------------
% Earlier notes (kept 2026-10-28; superseded by the draft above).
%
% Correctness estimation is a fundamental problem in machine learning
% evaluation and reliability assessment. It is the problem of predicting
% whether a model output matches the ground truth.
% In this work, we define code correctness as adherence to a given problem
% specification, which can be verified through test execution -- the standard
% measure used in code generation benchmarks and in practice.
% there are several approaches to correctness estimation of LLM generations,
% for natural language and code, with the best performing methods being based
% on internal representations
% But these methods are fragile under distribution shift
% Motivation: Steering agentic development (cheap signals for decisions within
% an agentic development loop)
% triage vs selection
% Comparable setup for introspective correctness prediction for code
% Nothing found shows probes degrading under benchmark or language shift for
% code. That gap is what our experiments fill, which is worth saying
% explicitly in the intro.
% Draft, overarching question (2026-10-24): We ask which execution-free
% signals of correctness survive benchmark and language shift, and what
% explains those that do. ...
% Übersichtsbild: Pipeline oder intervention
%
% From main.tex (2026-10-22 drafts):
% - Comparison of X execution-free correctness predictors from Y families, on
%   four (two?) models, five benchmarks in three programming languages, at
%   candidate level and at problem level.
% - Fragile under shift: internal representations have been used to estimate
%   the knowledge stored in LLMs, predict the risk of hallucination, predict
%   jailbreaking behaviours, and assess code correctness. However, a known
%   weakness of these approaches is their poor generalization behaviour across
%   tasks and datasets [OrgadTGRSKB25] -> focus on generalization.
% - 2 Questions: (1) Causal localization was proposed as a shift-robust
%   alternative for natural language [Huang2025Internal], is it beneficial for
%   code generation? (2) Across families, which degrades least, under which
%   type of shift (on which target)?
% - Draft contributions: comparison of six predictors from four families
%   (RQ1, RQ3, RQ4); adaptation of counterfactual simulation and value probing
%   to code, interventions at the last prompt token; analysis of where the
%   causal methods help (RQ2), incl. the withdrawn "the subspace encodes the
%   model's answer rather than the program's correctness".


\section{Related Work}
\label{sec:related}

\paragraph{Code correctness prediction}
Token probabilities form the most direct signal for execution-free code correctness prediction. They have been studied for calibration~\cite{Spiess2025Calib,Campos2026Multicalib}, but have been found to barely separate failing from passing programs~\cite{rathnasuriya2026overconfident}. Consistency methods compare multiple  generations for the same problem, some by executing them~\cite{valentin2026incoherence,sharma2025,he2026usingsemanticdistanceestimate}, others without execution~\cite{jain2023lightweight,shi2026codemorethantext,bouchard2026functional}. Probes and representation reading recover correctness from hidden states~\cite{bui2025openia,Vu2026Autoprobe,ribeiro2026repe,huang2025pttrust,lugoloobi2026llms}.
%\orange{Thomas?~\cite{klassert2026introspective}}. 
Comparisons across the above approaches only exist between output and consistency methods so far~\cite{bouchard2026functional}. %\orange{streichen/Relevanz? or set probes against sampling at line level~\cite{gros2025localized,huang2025pttrust}}; 
To our knowledge, no causal methods have been investigated.

\paragraph{Verdict correctness prediction}
LLMs are increasingly used as judges in software engineering~\cite{he2025judgese}, but
%and have been compared with human ratings on several SE tasks~\cite{wang2025replace}.}
 have been found to misjudge code correctness~\cite{crupi2025judge} and react to superficial edits such as renamed variables~\cite{moon2026cover}. %He et al.\cite{he2025judgese} survey their use in software engineering. 
Hence, {\it verdict correctness prediction} aims at predicting the correctness not of the code itself, but of the LLM's asssessment. 
Approaches have explored 
%The trustworthiness of a single verdict has been predicted using probes on the 
judges' hidden states~\cite{Radharapu2025CalibratingJudges}, reasoning structure~\cite{qi2026verdi}, or have escalated uncertain cases~\cite{jung2025trust}. 
%While internals help when a judge is comparing answers, they offer no advantage when scoring one answer at a time~\cite{mukherjee2026opening}. 
Furthermore, probes trained for correctness can track the model's own judgment instead~\cite{lu2026diagnosing}, which supports keeping the verdict and program correctness as separate targets. To our knowledge, however, no prior work studies both code correctness and verdict correctness, nor has verdict correctness been assessed under distribution shift. 

%%%% VIOLA? predicts the reliability of a code judge's verdict from its \orange{internals (Korrekt - Mukherjee?)}, or under distribution shift, and 

%Such probes have been trained on general judging data that includes coding items~\cite{Radharapu2025CalibratingJudges}; to our knowledge, no work predicts from its internals whether a model's verdict on its own code is right, or tests this under a change of benchmark or programming language.
%no prior work addresses both correctness and verdict correctness. 

\paragraph{Generalization under shift}
ML models in software engineering are well known to erode for cross-dataset and cross-project prediction~\cite{zimmermann2009cpdp,hosseini2019cpdp}, under language shift~\cite{he2026codecircuit,huang2025pttrust},
and under other realistic splits~\cite{ding2025primevul}.
This has also been found for recent LLM-based predictors ~\cite{OrgadTGRSKB25,Kramar2026Probes,Srey2026Factorised}, particularly for 
%often degrade outside their training distribution, across datasets and under production shifts, so their performance in distribution says little about their robustness
%Predictors trained on hidden states often degrade outside their training distribution: error detectors fail to transfer across datasets~\cite{OrgadTGRSKB25}, probes fail under realistic production shifts~\cite{Kramar2026Probes, Srey2026Factorised}.
%For code, %internal correctness signals have been evaluated under shift for one method family at a time: 
%correctness directions transfer poorly between fitting sources and benchmarks
execution-free code correctness prediction:  Probes based on hidden states have been found to lack robustness~\cite{ribeiro2026robustness}, and to fall behind under dataset shift~\cite{gros2025localized}.
%For code, correctness directions transfer poorly between fitting sources and benchmarks~\cite{ribeiro2026robustness}, attribution-graph features lose much of their signal across programming languages~\cite{he2026codecircuit}, %a hidden-state probe did not consistently beat token confidence after a transfer from Python to Java~\cite{huang2025pttrust}, 
%and probes that beat sampling in distribution fall behind it on a held-out dataset~\cite{gros2025localized}. %AUTOPROBE learns where to read in order to counter such fragility~\cite{Vu2026Autoprobe}. 
Each of these works studies one or two families of predictors on program correctness; we compare four families of predictors, evaluate not only code correctness but also verdict correctness, and include both benchmark and language shifts.


\paragraph{Causal localization}
%Distributed alignment search~\cite{GeigerWPIG24DAS} learns a subspace through which interchange interventions~\cite{Geiger2021CausalAbstractions} control a model's behavior; Huang et al.\cite{Huang2025Internal} use such subspaces to predict whether LLM outputs are correct and report robustness under shift on natural-language tasks; to our knowledge, this has not been applied to code. 
Huang et al.~\cite{Huang2025Internal} use causal localization 
%localize an LLM's internal state variables 
%decision variable \orange{decision variable finde ich komisch} 
with distributed alignment search~\cite{GeigerWPIG24DAS} to learn a subspace through which interchange interventions~\cite{Geiger2021CausalAbstractions} control the output, and use it for correctness prediction,  reporting strong robustness under shift on natural-language tasks. 
In general, whether a causal role makes a direction generalize better is open: some work finds that it does~\cite{marks2024geometry,mallen2024quirky}, while others find that causal efficacy and generality diverge~\cite{ying2026truthfulness}.
To this end, our work is the first to adapt and evaluate causal probing for code correctness and verdict estimation.
%this approach has not been applied to code. 

That an intervention changes the output does not by itself show that the model represents the variable there: patched subspaces can act through pathways the model does not normally use, and they need not be unique~\cite{makelov2024illusion,grant2025divergent,meloux2025everything}. We therefore compare the learned subspace with random subspaces of the same dimension % and \orange{test what it carries behaviorally 
(Section~\ref{sec:rq-subspace}). 





%Subspace interventions can activate dormant pathways\cite{makelov2024illusion,wu2024reply} and push representations off their natural distribution~\cite{grant2025divergent}, and alignments are not unique~\cite{meloux2025everything} or, without a linearity constraint, trivially perfect~\cite{sutter2025nonlinear}. 
%A subspace that changes the output when it is patched is not necessarily where the model represents the variable: the intervention can produce the effect through pathways the model does not normally use, or move the hidden state into regions that real inputs never reach~\cite{makelov2024illusion,grant2025divergent} (but see~\cite{wu2024reply}). Moreover, several different subspaces can fit equally well~\cite{meloux2025everything}, and with non-linear maps almost any network can be fitted to any hypothesis~\cite{sutter2025nonlinear}. We therefore keep the subspace linear and compare it with random subspaces of the same dimension.
%We therefore keep the subspace linear and compare it with random subspaces of the same dimension. Directions with a stronger causal role have been found to generalize better~\cite{marks2024geometry,mallen2024quirky}, but causal efficacy and generality can also diverge~\cite{ying2026truthfulness}, consistent with our finding that interchange accuracy does not indicate usefulness. Unlike subspace selection for probe transfer~\cite{yoo2026subspace}, our subspaces use no data from the target language.

\section{Problem Setting}
\label{sec:problem}
% Rewritten 2026-10-29 after Adrian's comments (define "correct" first; the
% judging pass is a definition, not a method; first shift paragraph was
% related work). Previous version kept as comments at the end.
We study execution-free correctness prediction for code that an LLM generates
and judges itself, from signals of the model's own generation process: no tests are
executed and no additional reasoning tokens are generated. Given a programming
problem $x_i$, a model $\mathcal{M}$ samples candidate solutions
$c_{ij} \sim \mathcal{M}(x_i)$, $j = 1, \dots, n_i$. Following common practice in code generation benchmarks, a candidate is
considered \emph{correct} if it passes all tests. We denote the program correctness label as $y^{\mathrm{prog}}_{ij} \in \{0,1\}$.
Tests are executed only to obtain these labels for training and evaluation,
never to make a prediction. Where a statement holds for any problem and
candidate, we write $(x, c)$ and drop the indices.

\subsection{The Model's Verdict} % \orange{Heißt später "confidence".}}
\label{sec:problem-verdict}
To obtain the model's own judgment, we prompt $\mathcal{M}$ with the problem description $x_i$ and the candidate solution $c_{ij}$
and ask whether the implementation is correct, explicitely demanding a single token answer, \texttt{Yes} or \texttt{No}. We define the model's \emph{verdict} $v(x, c)$ as the
answer $t \in \{\texttt{Yes}, \texttt{No}\}$ with the higher probability
$p_{\mathcal{M}}(t {\mid} x, c)$. We refer to this probability as the {\it verdict confidence} and use it as a baseline in our experiments. The same forward pass that defines the verdict -- the \emph{judging
pass} --
 also generates the signals for the other prediction methods (Section~\ref{sec:methods}). We extract those at the verdict token's position: 
%Because the prompt ends in the same fixed
%text for every input, the verdict is produced at one position, the last
%prompt token, and the methods differ only in what they read there: the
probabilities of the two answer tokens {\tt Yes/No}, the hidden states, or a subspace of
the hidden states identified through interventions. 
%EigenScore and LAT use their own inputs.
%,
%see Section~\ref{sec:methods}.

\subsection{Prediction Targets}
We evaluate predictors for two different targets: % that capture different notions of correctness.

\paragraph{Program correctness ($y^{\mathrm{prog}}$)}
The first target is to predict whether the candidate implementation actually satisfies the problem specification, as operationalized by the test suite. This target corresponds to prior work on code correctness
prediction~\cite{bui2025openia, ribeiro2026repe, valentin2026incoherence,
sharma2025, he2026usingsemanticdistanceestimate}, and 
prediction scores are used to decide which candidates need testing.
Note that predictors may differ from the model's verdict $v$ as defined in Section \ref{sec:problem-verdict}, since they can be based on other internal signals. 

\paragraph{Verdict correctness ($y^{\mathrm{verd}}$)}
The second target is whether the model's verdict (as defined in Section \ref{sec:problem-verdict}) is correct, that is, whether
the model accepts a correct candidate or rejects an incorrect one:
\begin{equation*}
y^{\mathrm{verd}} = \mathbf{1}\bigl[v(x, c) = \texttt{Yes} \Leftrightarrow
y^{\mathrm{prog}} = 1\bigr].
\end{equation*}
A predictor for this target estimates when the model's judgment can be
trusted and when to invoke more expensive verification, such as test
execution or a stronger model. It is related to uncertainty estimation for generated answers, which predicts whether a model's own answer is
correct~\cite{Kadavath2022LanguageM, Huang2025Internal}.

Note that these two targets serve different purposes: program correctness concerns the
candidate itself, whereas verdict correctness concerns the reliability of the model's
judgment. When a model correctly rejects an incorrect candidate, for example,
its verdict is correct while the program is not. We evaluate both targets with the
same judging pass to investigate which internal signals support each prediction target.

\subsection{Distribution Shift}
\label{sec:problem-shift}
%A central challenge for execution-free correctness prediction is generalization beyond the training distribution, since a predictor may exploit regularities of its training data that do not hold
%elsewhere~\cite{OrgadTGRSKB25, Srey2026Factorised}. 
We study two
generalization scenarios, in which the evaluation data differ from the
training data in one controlled way.

In \emph{benchmark shift}, we train a predictor on one benchmark,
LiveCodeBench~\cite{Jain2025Livecodebench} or
BigCodeBench~\cite{Zhuo2025BCB}, and evaluate it on the other. This simulates
a predictor fitted on one kind of task and deployed on another:
LiveCodeBench consists of algorithmic programming challenges, BigCodeBench of
practical tasks that call libraries and APIs.

In \emph{language shift}, we train a predictor on candidates in one
programming language and evaluate it on candidates in another, for the same
problems in Python, C++ and Java. The problems are split once for all three
languages, so no problem is seen in training in one language and in
evaluation in another. 
%The shift thus changes the language of the code while
%the tasks stay comparable.

% ---------------------------------------------------------------------------
% Previous version (until 2026-10-28), condensed. The judging-prompt listing
% moved to appendix.tex (label app-listing:prompt).
%
% We study execution-free code correctness prediction as a means of steering
% agentic software development without repeatedly invoking tests or other
% verification tools. In particular, we investigate whether these signals
% generalize across benchmarks and programming languages, rather than relying
% on correlations specific to the training distribution. ... We construct a
% dataset D = {(x_i, c_ij, y_ij^prog)} over N problems, using test execution
% only to establish labels, not to make predictions. ... Such estimates could
% help prioritize candidates for verification or guide decisions within an
% agentic loop.
%
% Shared Judging Pass: We evaluate all methods on a common judging prompt ...
% The methods differ only in what they read there: the probabilities of the
% two answer tokens, the hidden states, or a subspace of the hidden states
% identified through interventions (Section III).
%
% Program correctness: ... A predictor for this target estimates whether the
% candidate is correct, regardless of the model's own assessment. Such a score
% can support candidate selection and help determine which generated programs
% require further verification.
% Verdict correctness: ... This target is related to uncertainty estimation for
% model-generated answers ... [Kadavath2022LanguageM, Huang2025Internal, TODO].
%
% Distribution Shift, first paragraph: A central challenge for execution-free
% correctness prediction is generalization beyond the training distribution.
% Probes on hidden states often fail to generalize across datasets
% [OrgadTGRSKB25, Levinstein2025LieDetector] and under production shifts
% [Kramar2026Probes]; for code, probe-based correctness assessors lose
% accuracy when applied to a model other than the one they were trained on
% [Vu2026Autoprobe]. A predictor may exploit regularities specific to the
% training problems, benchmark, or programming language rather than signals
% that reliably indicate correctness. Performance on held-out problems from
% the same benchmark therefore does not establish that a predictor will remain
% effective under distribution shift [Srey2026Factorised].
% In benchmark shift, the source and target domains consist of different
% benchmarks. ... In language shift, ... using a shared partition of the
% underlying problems to prevent overlap between training and evaluation
% problems. This setting tests whether correctness signals transfer across
% programming languages, where code representations and language-specific
% implementation details differ while the underlying problems correspond.


\section{Methods}
\label{sec:methods}
% Rewritten 2026-10-29 after Adrian's comments (base/source introduced first,
% training pairs explained for correctness first, intuition of counterfactual
% simulation before the formulas, LAT labels and orientation, Table I single
% column). Previous version: git show e509b60:paper/sections/methods.tex
We compare six predictors from four families, grouped by the signal they read
from the model: (1) its output token probabilities, (2) the consistency across
sampled generations, (3) internal features of the hidden states, and (4)
internal causal features, a subspace of the hidden states identified through counterfactual interventions. The first three families are established approaches to
correctness prediction. The fourth was proposed for natural-language question
answering by Huang et al.~\cite{Huang2025Internal}, who report that it is more robust under distribution shift. We adapt this approach to code generation %resulting in a new method for code correctness prediction 
and describe it in more detail.


\subsection{Output Token Probabilities (Confidence)}
\label{sec:methods-output}
The simplest predictor uses the probability the model assigns to an answer in the
judging pass. It requires neither hidden states nor labeled training examples,
and we refer to it as \emph{confidence}. For program correctness, we use the
probability of \texttt{Yes}, $p_{\mathcal{M}}(\texttt{Yes} \mid x, c)$,
following 
%the P(True) approach of 
Kadavath et
al.~\cite{Kadavath2022LanguageM}. For verdict correctness, we use the
probability of the emitted verdict, $p_{\mathcal{M}}(v(x, c) \mid x, c)$.
%, regardless of the generated answer. 
A model that is uncertain of its verdict should be
wrong more often~\cite{Hendrycks2017Baseline}. Both come from a softmax over the full vocabulary, whose
temperature is the only free parameter (Section~\ref{sec:setup-training}).

\subsection{Consistency (EigenScore)}
\label{sec:methods-consistency}
Consistency based methods do not judge a single candidate but asks how consistent
the model is across several generations for the same problem: a model that knows
how to solve a problem should produce similar solutions, while diverging
samples indicate uncertainty. 
%Most consistency measures for code execute the
% generated programs (Section~\ref{sec:related}) to determine similarity. 
We use EigenScore~\cite{Chen24Eigenscore},
which measures consistency in the hidden states. For a
problem with $m$ sampled candidates, it takes the hidden state of the last
layer at the last token of each candidate's generation and computes the
log-determinant of their regularized covariance matrix, normalized by $m$: the
more the states spread out, the higher this {\it Eigenscore}. We use the negated
EigenScore as an indicatore that 
%the score that the problem is solved. 
instances from the candidate set solve the problem. 
EigenScore is an unsupervised method, requiring no correctness labels. Unlike all other methods, it scores a problem rather than a candidate and does not rank the candidates of one problem. We compare it with the other methods at problem level (Section~\ref{sec:rq-problem}). 
Its cost is $m$ generations per problem, which come for free wherever multiple candidates are sampled anyway.

\subsection{Internal Features}
\label{sec:methods-internal} 
The third family reads correctness from the hidden states. We extract the
hidden state $h_\ell(x, c) \in \mathbb{R}^d$ after a selected layer $\ell$ at the
last token of the judging prompt, the position from which the verdict is
generated. $h_\ell(x, c)$ serves as input for two probing methods:
%Different layers have been found to encode different aspects of a computation~\cite{JinYHZWH0MMDYDZ25}, so the choice of layer can substantially affect predictive performance. We fit the following two methods at every layer and choose the best performing layer \orange{in distribution (d.h. auf dem Testset?)} (Section~\ref{sec:setup-training}).

\paragraph{Correctness probe}
We use linear probes~\cite{Alain2016Probes}, i.e. logistic regressors
trained on $h_\ell$ with correctness labels~\cite{OrgadTGRSKB25,
bui2025openia, Vu2026Autoprobe}: $y^{\mathrm{prog}}$ for program correctness and $y^{\mathrm{verd}}$ for verdict correctness. A probe's predicted probability of the positive class serves as the score. 
Since the classifier is linear, it captures only linearly accessible information~\cite{park2024linear}, like the other hidden states methods.

\paragraph{Linear Artificial Tomography (LAT)}
% Reported LAT = stimulus "lat" of run_lat_sweep.py (Ribeiro et al.'s
% task+code stimulus under its meta-prompt), all models, folds and benchmarks.
% The judging-prompt variant (stimulus "verdict", method lat_verdict in the
% artifact) exists for gpt-oss fold 0 on the language sets only: similar in
% distribution, lower under shift (best ID layer, not cross-fitted; checked
% 2026-10-29).
LAT~\cite{ribeiro2026repe, Zou2023RepE} identifies a correctness direction:
%without
%fitting a classifier, but uses labels to construct and orient this direction. 
For each training problem with both a passing and a failing candidate, it computes the difference
of their hidden states and applies principal component analysis (PCA) to these difference vectors. 
The first principal component defines the direction, whose sign is chosen so that passing candidates tend to have a higher projection than failing ones. 
%The labels also fix its orientation: the direction is signed so that, in most training pairs, the passing candidate projects higher. 
A candidate's score is its projection onto this direction.
In the original method, hidden states are extracted using a dedicated stimulus prompt that presents the problem and candidate~\cite{ribeiro2026repe}. In a preliminary comparison, applying the same procedure to the judging pass did not improve performance, so we use the original prompt for fairness. LAT was designed to predict program correctness, and we evaluate it for this target only.


\subsection{Internal Causal Features}
\label{sec:methods-causal}
These methods follow the tradition of causal interpretability, which identifies a
causal representation by intervening on it~\cite{Geiger2021CausalAbstractions}. 
Here, these approaches try to localize specific features within hidden states that causally determine a model's verdict rather than merely correlating with correctness. If a model judges code reliably when prompted, it presumably does so through a systematic internal
computation. 
We ask whether this computation can be localized 
%on candidates the model judges correctly, 
and whether it helps predicting correctness on new inputs.

\paragraph{Localizing the verdict}
We assume that the model computes an internal assessment of whether a given
candidate is correct, a variable $V \in \{\texttt{Yes}, \texttt{No}\}$ that
determines its verdict and is encoded in the hidden state at the last prompt
token. To localize it, we use Distributed Alignment Search (DAS)~\cite{GeigerWPIG24DAS}, which can be understood by analogy to debugging: 
Run a program on two inputs, pause both at a breakpoint, copy a variable's value from the second run into the first, and continue: if the first run now produces the expected output of the second, the variable {\it controls} the output. 
The run that is modified is called the \emph{base} $b = (x_b, c_b)$, the run whose value is copied the \emph{source} $s = (x_s, c_s)$. 
Note that both are usually sampled from different problems (Figure~\ref{fig:intervention}). 

% Source: figures/fig_intervention.tex (standalone TikZ, 2026-10-29).
\begin{figure}[t]
\centering
\includegraphics{generated/fig_intervention}
\caption{(a) The judging pass: the model reads the problem, the candidate and
the question and answers \texttt{Yes} or \texttt{No} at the last prompt token.
The confidence baseline reads answer probabilities, the probe the hidden state
$h_\ell$ of one layer, the value probe its projection onto the learned
subspace. 
(b) An interchange intervention: the subspace coordinates of a
source example's hidden state replace those of a base example, and the base's forward pass continues. 
DAS learns a subspace so that the base then answers like the source.
Counterfactual simulation scores a candidate by making it the source for $B = 16$ fixed bases from the training set. A source state that firmly encodes its verdict should enforce it on the bases, one on which the model is unsure should not. The more consistently the bases adopt the candidate's verdict, the more likely it is correct. 
%Counterfactual simulation repeats the intervention with $B {=} 16$
%fixed bases and averages the probability that they answer like the source. \orange{Klarmachen warum/dass damit eine Prognose über das Label der *Source* getroffen wird?}
}
\label{fig:intervention}
\end{figure}

In an LLM there are no named variables whose values could be switched as in the debugging analogy, and a concept such as ``this program is correct'' is typically spread across many dimensions of $h_\ell$. 
DAS therefore \emph{learns} where to look: a $k$-dimensional subspace, given by $R \in \mathbb{R}^{d \times k}$ with orthonormal columns. 
The intervention replaces the  component of the base state's representation within this subspace with the corresponding component from the source
% base's coordinates inside the subspace with the source's 
and leaves the rest of the base's state untouched:
\begin{equation}
  h_\ell^{\,b\leftarrow s} \;=\; h_\ell(b) + R R^{\top}\bigl(h_\ell(s) - h_\ell(b)\bigr).
\end{equation}
The forward pass of $b$ then continues through all later layers from
$h_\ell^{\,b\leftarrow s}$ instead of $h_\ell(b)$. Such a replacement is
called an interchange intervention~\cite{Geiger2021CausalAbstractions}.

\paragraph{Training}
DAS is trained only on candidates the model judges correctly ($y^{\mathrm{verd}}{=}1$). 
On these, the verdict coincides with program correctness $y^{\mathrm{prog}}$: \texttt{Yes} for the passing
candidates, \texttt{No} for the failing ones. 
A training pair combines two such candidates as base and source, from the same or different problems and with the same or different verdicts. 
With the LLM frozen, $R$ is trained so that the intervened base outputs the source's verdict: Where the two verdicts differ, the intervention must flip the base's answer; where they agree, it
must leave it unchanged. The share of held-out interventions that produce the
intended verdict is the interchange intervention accuracy (IIA)~\cite{GeigerWPIG24DAS}. 
It shows whether the subspace controls the verdict, not how well the verdict predicts correctness. We learn one subspace per layer and use it for both targets.

\paragraph{Counterfactual Simulation}
Counterfactual simulation~\cite{Huang2025Internal} asks whether a (source) candidate's
state in the subspace is strong enough to impose its assessment on other
(base) examples. The intuition here is that a source candidate whose state clearly encodes a verdict -- \texttt{Yes} or \texttt{No} --
should do so, while one on which the model is unsure should not. 
We fix a set $\mathcal{B}$ of $B {=} 16$ correctly judged training candidates as bases. 
The candidate to be scored is the source $s$. 
Let $p_{b \leftarrow s}(t)$ be the probability of answer $t$ in the intervened
forward pass of base $b$. Averaged over the bases, we calculate two scores:
\begin{align}
\mathrm{cs}_{\mathrm{prog}}(s) &= \frac{1}{B}\sum_{b \in \mathcal{B}}
p_{b \leftarrow s}(\texttt{Yes}), \\
\mathrm{cs}_{\mathrm{verd}}(s) &= \frac{1}{B}\sum_{b \in \mathcal{B}}
p_{b \leftarrow s}\bigl(v(s)\bigr).
\end{align}
The first measures how strongly the candidate's state pushes the bases' verdicts toward
\texttt{Yes} and predicts program correctness. 
The second measures how consistently the bases adopt the candidate's own verdict and predicts verdict
correctness. Note that counterfactual simulation requires $B$ additional forward passes per candidate.

\paragraph{Value Probe}
The value probe~\cite{Huang2025Internal} utilizes the subspace $R$ for probing instead of
intervening on it, so it needs no forward pass beyond the judging pass. Like the correctness probe (Section \ref{sec:methods-internal}), it uses logistic regression, but on
the $k$-dimensional projection $R^{\top} h_\ell(x, c)$ of the candidate's hidden state onto subspace $R$. %trained to predict the
% model's verdict on the correctly judged training candidates. On these, the
% verdict equals the program label, so the probe learns which values of the
% subspace go with correct and with incorrect programs when the model judges
% right. Its probability of \texttt{Yes} scores program correctness; its
% probability of the candidate's actual verdict scores verdict correctness. It
% needs no forward pass beyond the judging pass.
This probe is trained differently for the two targets:
For the verdict target, it learns which verdict a subspace state encodes. It is trained on the same candidates as DAS, those the model judges correctly,
to predict the model's verdict, which on these candidates equals $y^{\mathrm{prog}}$: \texttt{Yes} for passing and \texttt{No} for failing programs. Following Huang et al., the
score is the probe's confidence in its own prediction, $|p(\texttt{Yes}) - \tfrac12|$. As in counterfactual simulation, a state that clearly reads as \texttt{Yes} or \texttt{No} indicates certainty and a probably correct verdict.
For the program target, we train the probe on all training candidates to predict $y^{\mathrm{prog}}$, and its probability of \texttt{Yes} is the score. 
%It is a correctness probe restricted to the subspace, which lets us compare the learned subspace with random subspaces of the same dimension (Section~\ref{sec:rq-subspace}). 



% \paragraph{Adaptation to code}
% Huang et al.~\cite{Huang2025Internal} study tasks with short prompts and strongly constrained answers in natural language. In our setting, base and source contain different problems and programs, so they differ in length, and positions do not correspond. 
% The judging prompt provides one position that does: the last prompt token is identical in
% every input, and the verdict is predicted from its hidden state. We therefore
% intervene there only. At an intermediate layer, later layers can still read
% the program tokens through attention, so the intervention controls the
% verdict only insofar as the relevant information has already been gathered at
% this position. The share of \texttt{Yes} and \texttt{No} verdicts in the
% training pairs is set by the model and the benchmark; we do not rebalance it. \orange{Könnte man meiner Meinung nach stark kürzen.}



% \begin{table}[t]
% \centering\footnotesize
% \setlength{\tabcolsep}{2.5pt}
% \caption{What the methods require. \emph{Labels}: supervision used for
% fitting; every method also uses labeled in-distribution data to choose its
% layer or temperature. Correctly judged candidates are identified with the
% execution labels. \emph{Extra cost}: beyond the judging pass.}
% \label{tab:method-requirements}
% \begin{tabular}{lllll}
% \toprule
% Method & Labels & Fitting & Extra cost & Scores \\
% \midrule
% Confidence & -- & -- & -- & candidate \\
% EigenScore & -- & -- & $m$ samples & problem \\
% \midrule
% Probe & pass/fail & log.\ reg. & -- & candidate \\
% LAT$^\dagger$ & pass/fail pairs & PCA & own prompt & candidate \\
% \midrule
% Counterf.\ sim. & correctly judged & DAS & $B$ passes & candidate \\
% Value probe & correctly judged$^\ddagger$ & DAS + log.\ reg. & -- & candidate \\
% \bottomrule
% \end{tabular}
% \par\smallskip
% \parbox{\linewidth}{\scriptsize $^\dagger$Program target only; reads its own
% stimulus prompt instead of the judging pass. $^\ddagger$The subspace is learned
% on correctly judged candidates; the regression on it uses the correctly judged
% candidates with their pass/fail labels for the verdict target and all candidates for the program target. Fitting is per layer for all hidden-state methods.}
% \end{table}



\section{Study Design}
% RQ subsection removed 2026-10-29: the RQs are stated in the introduction.
We evaluate the methods of Section~\ref{sec:methods} on both targets of
Section~\ref{sec:problem}, verdict and code correctness, in distribution and under benchmark and language
shift. RQ1--RQ3 judge single candidates, RQ4 whole problems.
% \subsection{Research Questions}
% \label{sec:setup-rqs}
% We ask which execution-free signals of correctness survive benchmark and
% language shift, whether causal localization makes them more robust, and what explains the signals that survive. To answer this, we compare
% the methods of Section~\ref{sec:methods} on the verdict and program target of
% Section~\ref{sec:problem}, in distribution and under benchmark and language
% shift (Section~\ref{sec:problem-shift}). RQ1 compares the methods on the
% verdict target, RQ2 examines what the causal methods' predictions rest on,
% RQ3 asks whether the same readouts predict program correctness, and RQ4 moves
% from candidates to problems.

% \smallskip\noindent\textbf{RQ1:} \emph{How well can execution-free methods
% tell whether the model's own verdict can be trusted, in distribution and under
% shift?} A predictor for the verdict target tells whether the model's
% judgment can be trusted or must be checked by execution. %It is the target the causal methods were designed for~\cite{Huang2025Internal}.
% We train the probe and the two causal methods on one benchmark or language
% and evaluate them, together with confidence, there and on the others. 

% \smallskip\noindent\textbf{RQ2:} \emph{Which part of the causal methods
% carries their signal, and what does the learned subspace encode?} Each causal
% method combines a subspace learned at one layer, an intervention through it,
% and a readout of the model's answer token. We separate these parts, replace
% the learned subspace by random ones of the same dimension, and ask whether the
% subspace encodes the model's verdict or the correctness of the program. %We examine this where RQ1 finds the causal methods ahead of the non-causal ones, with the other cases as the contrast.

% \smallskip\noindent\textbf{RQ3:} \emph{Do the same readouts tell whether the
% program is correct?} A predictor for the program target tells whether a
% candidate can be accepted without running it. This target coincides with the
% verdict target only where the model judges correctly, and it is the target
% most prior work addresses. A subspace that encodes the model's verdict can
% predict program correctness only as far as the verdict does. If RQ2 finds
% that the subspace encodes the verdict, the causal methods should offer no
% advantage on this target. We test this with the same readouts
% and evaluate every single-pass method, including LAT, under the same shifts.

% \smallskip\noindent\textbf{RQ4:} \emph{Which problems will the model solve?
% A single pass against agreement across samples.} A problem-level prediction
% tells where to spend testing effort, or whether to sample more candidates.
% When several candidates are sampled for a problem, their agreement is a
% further signal that needs no labels and no training, so unlike the trained
% methods of RQ1 and RQ3 it cannot lose under shift. We compare EigenScore~\cite{Chen24Eigenscore}
% with the single-pass methods at problem level and at equal sampling cost, and
% with a probe that reads the model before it generates anything.


\subsection{Models and Data}
\label{sec:setup-data}

\paragraph{Models}
We study two open-weight models, \texttt{Qwen3-Coder-30B-A3B-Instruct} and
OpenAI \texttt{gpt-oss-20b}. Each model generates candidate programs and judges its own candidates, so generation, verdicts and hidden states come from the same model. 
Qwen3-Coder answers the judging prompt directly. 
Since gpt-oss is a reasoning model, we prefill its final channel so that it also answers directly. 
Preliminary experiments showed that reasoning would cost about 1{,}100 generated tokens per candidate (median at medium effort) instead of a single forward pass, and its varying
length would shift the position at which the methods read the hidden state.
Judging without reasoning costs gpt-oss some accuracy (Section~\ref{sec:threats}).

\paragraph{Benchmarks}
We construct our datasets from LiveCodeBench~\cite{Jain2025Livecodebench} (\texttt{release\_v6}) and BigCodeBench~\cite{Zhuo2025BCB}, both written in Python. 
%and HumanEval~\cite{chen2021humaneval}, all in Python. 
To evaluate  language shift we use 454 LiveCodeBench problems released between October 2024 and May 2025. We use the original Python versions, together with corresponding C++ versions from the LiveCodeBench-CPP dataset (\texttt{v6\_2408\_2505})\footnote{\url{https://huggingface.co/datasets/nvidia/LiveCodeBench-CPP}} and Java versions sourced from the original AtCoder and LeetCode platforms. We refer to these datasets as LCB-Python, LCB-C++, and LCB-Java, respectively. LCB-Python is the subset of our LCB data on these problems.



\paragraph{Candidates and labels}
For each problem, each model samples $n = 10$ candidates at temperature $t=1$. We discard generations from which no code candidate can be extracted, candidates whose judging prompts exceed 2048 tokens, and duplicate generations.
% TODO hier?
EigenScore compares $m$ non-empty candidates per problem: $m = 5$ on LCB and
the language sets, where some generations are empty, and all ten on BCB,
where almost none are. Problems with fewer non-empty candidates are left out
of the problem-level comparison for every method (93--100\% of the problems
are kept for Qwen3-Coder, 76--94\% for gpt-oss).

We obtain program-correctness labels by executing each candidate against the respective benchmark's test suite. To obtain verdict-correctness labels, we prompt the generating model to judge each candidate and compare its verdict with the execution-based correctness label
(Section~\ref{sec:problem}).

\paragraph{Splits}
We split each benchmark by problem into training, validation and test sets (60/20/20), so all candidates of a problem fall into the same split. 
For the language sets, with fewer than 90 test problems each, we evaluate on validation and test together. 
We repeat every experiment on three partitions of the
problems with disjoint test sets and report the mean over the partitions.


Table~\ref{tab:datasets} summarizes the evaluation sets of the first
partition and reports the fraction of candidates that pass the tests, the
fraction accepted by the model, and the balanced accuracy of the model's
verdict against program correctness. The models' balanced accuracy ranges from
57.2\% to 74.4\%, so the models misjudge many of their own candidates, and
knowing which verdicts to trust matters.

\begin{table}[t]
\centering
\caption{Evaluation sets of the first partition. \emph{Problems}: number of distinct problems; \emph{Cand.}: number
of scored candidates; \emph{Pass}: percentage of candidates that pass the tests; \emph{Yes}: percentage of candidates accepted by the model;
\emph{Judge acc.}: balanced accuracy of the model's verdict against program correctness (mean of the acceptance rate on passing and the rejection rate on failing candidates). }
\label{tab:datasets}
\begin{tabular}{l rrrrr}
\toprule
Benchmark & Problems & Cand. & Pass & Yes & Judge acc. \\
\midrule
\multicolumn{6}{l}{\emph{Qwen3-Coder-30B}} \\
LiveCodeBench & 208 & 1646 & 55.3 & 58.3 & 74.4 \\
BigCodeBench  & 456 & 3240 & 40.2 & 71.7 & 58.3 \\
LCB-Python    & 178 & 1423 & 41.7 & 44.1 & 73.8 \\
LCB-C++       & 175 & 1448 & 24.0 & 32.5 & 63.8 \\
LCB-Java      & 179 & 1486 & 37.5 & 31.2 & 64.4 \\
\midrule
\multicolumn{6}{l}{\emph{gpt-oss-20b}} \\
LiveCodeBench & 189 & 1588 & 80.4 & 44.6 & 62.8 \\
BigCodeBench  & 456 & 4533 & 45.8 & 42.6 & 61.7 \\
LCB-Python    & 150 & 1140 & 76.8 & 37.6 & 62.2 \\
LCB-C++       & 147 & 1116 & 68.4 & 35.1 & 57.2 \\
LCB-Java      & 141 & 1069 & 59.8 & 38.6 & 65.4 \\
\bottomrule
\end{tabular}
\end{table}

\subsection{Training and Selection}
\label{sec:setup-training}
% For the replication package, not the paper (checked 2026-10-24):
% - probe: LogisticRegression C=0.05 (src/probe.py); --sample-counts 512 =
%   per class (run_probe_sweep.py).
% - LAT: src/lat.py build_pairs, pairs_per_problem 1; stimulus prompt of
%   ../ribeiro_correctness.
% - DAS: --dims 6 --lrs 1e-4 --sample-counts 1024 --epochs 2
%   --early-stop-patience 6 (window 250 steps, min delta 0.01, on training
%   IIA; src/das_training.py). No run entered epoch 2 (Viola). Pairs: all
%   ordered pairs of the first 1024 judge-correct training candidates, base !=
%   source, same or different problem and verdict, shuffled, capped at 20,480
%   (causal_data_utils.load_intervention_data); target = source verdict; no
%   label balancing. IIA monitored on training bases x validation sources.
% - Counterfactual simulation: --num-verified-eval 16, --eval-bases head = the
%   first 16 judge-correct training candidates in file order (about two
%   problems, one verdict; DAS_EVAL_VALIDITY.md).
% - Value probe: LogisticRegression C=0.05 on R^T h of 1024 judge-correct
%   training candidates (run_value_probe_sweep.py).
% - Confidence: temperature grid differs by task (Qwen LCB and Python-control
%   verdict 0.5-4.5; LCB/BCB program 1.0, 2.0; else 1.0-2.0).
% - Selection: crossfit_layers.py (5 folds, 20 repeats);
%   score_layer.restrict_to_shared (candidates per partition).
% - Layer numbering: probe/LAT by hidden-state index (L = block L-1), DAS by
%   block (FINDINGS §10).
\paragraph{Fitting}
The probe is a logistic regression on the hidden state of one layer, trained
on 512 candidates per class. LAT uses the stimulus prompt of Ribeiro et
al.~\cite{ribeiro2026repe} and is fitted on one pair of a passing and a
failing candidate per training problem that has both. DAS learns a
$k {=} 6$-dimensional subspace per layer, a size chosen in a preliminary sweep
and in line with evidence that correctness is encoded in few
dimensions~\cite{cho2026confidencemanifold}. It is trained on pairs of 1024
correctly judged training candidates, in at most one pass over 20,480 pairs
with early stopping on the training IIA. The value probe is a logistic
regression on the projections of the same training candidates onto the subspace.
Counterfactual simulation uses a fixed set of $B {=} 16$ correctly judged
training candidates as bases. The generation probes of RQ3 apply the probe to
the hidden states of the model's own generation, replayed token by token, at
the last token of the code or of the prompt; they are trained on the same
candidates and program labels as the probe. Detailed settings can be found in the replication package\footnote{\url{https://anonymous.4open.science/r/FORGE-PCCuDS}}.

\paragraph{Selection}
Each method comes with one free parameter: the layer, or --
for confidence -- the temperature. The probe and LAT are fitted on every layer, while DAS
is trained at a subset (every second layer of L33--L43 for Qwen3-Coder and of L11--L23 for gpt-oss), since each subspace takes several GPU hours.
Layers are chosen in distribution, via cross-fitting over problems.
The selection is made on four fifths of the evaluation problems and read on the
remaining fifth, over 20 random splits. Under shift, the choice made on all
in-distribution rows is read on the shifted set, so no shifted row enters
selection. %\orange{The best layer on the shifted rows, which no selection rule can find in advance, appears only as a reference (Figure~\ref{fig:verdict-profile})}.
%We do not select on a separate validation split: on the language sets it is part of the evaluation set, which doubles the number of evaluation problems, and on LCB and BCB it stays unused so that one rule holds for all benchmarks; apart from measuring IIA it plays no role.

\subsection{Evaluation Protocol}
\label{sec:setup-protocol}
% Within-problem AUROC: make_latex_results.within_problem, all (positive,
% negative) pairs of the same problem pooled, ties 1/2; ID per held-out fold
% of the cross-fit, pair-weighted, mean over repeats (crossfit_layers.py).
\paragraph{Metrics}

At candidate level (RQ1--RQ3) we report the area under the receiver operating characteristic curve (AUROC)
over all candidates in an evaluation set, as is common for correctness prediction~\cite{he2026usingsemanticdistanceestimate,he2026codecircuit,bouchard2026functional,shi2026codemorethantext}, using either verdict correctness or program correctness as the positive label. AUROC measures how well a score ranks positive examples above negative ones, with 0.5 corresponding to chance and 1.0 to perfect ranking.
Because candidates from different problems are pooled, this metric reflects both the ability to distinguish candidates within a problem and the ability to rank problems by how often the model produces positive outcomes. Where this distinction matters, we additionally report within-problem AUROC, computed over pairs of candidates from the same problem with different labels. It equals the pairwise accuracy of~\cite{ribeiro2026robustness} and is on the same scale as the pooled AUROC.
The gap between the two is the part of the pooled AUROC that comes from ranking problems.
At problem level (RQ4), a problem is
considered as solved if at least half of its candidates pass, and the
single-pass methods are averaged over the candidates that EigenScore uses.

\paragraph{Significance}

We compare methods using a paired, problem-clustered bootstrap with 2,000 resamples. Each resample draws whole problems and evaluates all methods on the same examples, accounting for within-problem dependence while preserving paired comparisons. For language-shift results, we average differences over the six language pairs and three partitions, using shared resamples across language pairs because they share problems. We combine the bootstrap standard error (SE) with between-partition variability to account for variation due to training and layer selection. We consider a difference significant if $|\Delta| > 1.96\,\mathrm{SE}$ and report $\Delta \pm 1.96\,\mathrm{SE}$.



% Methods are compared on identical candidates with a paired, problem-clustered bootstrap: each of 2000 draws resamples problems and
% evaluates every compared method on the same draw. Resampling whole problems
% accounts for candidates of one problem not being independent; evaluating all
% methods on the same draw makes the comparison paired. On the language sets,
% we compare means over the six language pairs, averaged over the three
% partitions. The standard error (SE) of such a difference combines the
% bootstrap standard error with the spread of the difference between
% partitions, which reflects training and layer selection; the bootstrap alone,
% with the chosen layer fixed, does not capture this. A difference is
% considered significant if $|\Delta| > 1.96\,\mathrm{SE}$, and we report
% $\Delta \pm 1.96\,\mathrm{SE}$. Where only one partition exists (LCB, BCB and
% RQ4), a difference is considered significant if its 95\% bootstrap interval
% excludes zero. Results for individual language pairs are descriptive.



\section{Experiments}

\subsection{RQ1: How well can execution-free methods tell whether the model's own verdict can be trusted? }
\label{sec:rq-verdict}

% Generated by scripts/paper/make_result_tables.py -- do not edit.
\begin{table}[t]
\centering\footnotesize
\setlength{\tabcolsep}{3.2pt}
\caption{AUROC for verdict correctness (mean over 3 partitions). \emph{LCB}, \emph{BCB}: trained and evaluated on that benchmark; $\to$: evaluated on the other one. Language shift: \emph{ID} is the mean over the three training languages, \emph{Shift} the mean over the six OOD train$\to$evaluation pairs (Table~\ref{tab:language-detail-verdict}), $\Delta$ their difference. %In-distribution cells are cross-fitted over problems; shifted cells use the layer chosen in distribution. 
Best per column and LLM in bold.}
\label{tab:verdict-summary}
\begin{tabular}{l rrrr rrr}
\toprule
 & \multicolumn{4}{c}{Benchmark shift} & \multicolumn{3}{c}{Language shift} \\
\cmidrule(lr){2-5}\cmidrule(lr){6-8}
Method & LCB & $\to$BCB & BCB & $\to$LCB & ID & Shift & $\Delta$ \\
\midrule
\multicolumn{8}{l}{\emph{Qwen3-Coder-30B}} \\
Confidence & .672 & \textbf{.552} & .555 & .669 & .644 & .643 & \textbf{$-$.001} \\
Probe & .699 & .480 & \textbf{.639} & .516 & \textbf{.761} & .668 & $-$.093 \\
\addlinespace[2pt]
Counterf.\ sim. & .692 & .453 & .560 & .680 & .759 & \textbf{.728} & $-$.031 \\
Value probe & \textbf{.740} & .485 & .579 & \textbf{.695} & .758 & .728 & $-$.031 \\
\midrule
\multicolumn{8}{l}{\emph{gpt-oss-20b}} \\
Confidence & .569 & .598 & .597 & .544 & .542 & .536 & \textbf{$-$.005} \\
Probe & \textbf{.829} & .478 & .507 & .476 & .774 & .668 & $-$.106 \\
\addlinespace[2pt]
Counterf.\ sim. & .779 & .572 & \textbf{.601} & .556 & \textbf{.808} & \textbf{.691} & $-$.116 \\
Value probe & .695 & \textbf{.599} & .600 & \textbf{.605} & .713 & .646 & $-$.067 \\
\bottomrule
\end{tabular}
\end{table}

% Generated by scripts/paper/make_result_tables.py -- do not edit.
\begin{table*}[t]
\centering\footnotesize
\setlength{\tabcolsep}{1.5pt}
\caption{AUROC for verdict correctness, in distribution (ID) and under language shift per training language (LCB problems; ID cross-fitted, shifted columns at the layer chosen in distribution). Mean $\pm$ SD over 3 partitions of the problems. Best per column in bold; underlined: within one SD of the best. Non-causal methods above the rule, causal below.}
\label{tab:language-detail-verdict}
\begin{tabular*}{\textwidth}{@{\extracolsep{\fill}}ll rrr rrr rrr@{}}
\toprule
 & & \multicolumn{3}{c}{Trained on Py} & \multicolumn{3}{c}{Trained on C++} & \multicolumn{3}{c}{Trained on Java} \\
\cmidrule(lr){3-5}\cmidrule(lr){6-8}\cmidrule(lr){9-11}
Model & Method & ID & $\to$C++ & $\to$Java & ID & $\to$Py & $\to$Java & ID & $\to$Py & $\to$C++ \\
\midrule
\multirow{4}{*}{Qwen3-Coder-30B} & Confidence & .608\,{\scriptsize$\pm$.028} & .637\,{\scriptsize$\pm$.052} & .687\,{\scriptsize$\pm$.011} & .637\,{\scriptsize$\pm$.054} & .611\,{\scriptsize$\pm$.026} & .688\,{\scriptsize$\pm$.013} & .689\,{\scriptsize$\pm$.010} & .605\,{\scriptsize$\pm$.027} & .630\,{\scriptsize$\pm$.053} \\
 & Probe & .670\,{\scriptsize$\pm$.011} & \underline{.628}\,{\scriptsize$\pm$.103} & .636\,{\scriptsize$\pm$.059} & \textbf{.874}\,{\scriptsize$\pm$.010} & .636\,{\scriptsize$\pm$.062} & .701\,{\scriptsize$\pm$.009} & \underline{.740}\,{\scriptsize$\pm$.044} & .670\,{\scriptsize$\pm$.026} & .739\,{\scriptsize$\pm$.020} \\
\cmidrule(lr){3-5}\cmidrule(lr){6-8}\cmidrule(lr){9-11}
 & Counterf.\ sim. & \underline{.671}\,{\scriptsize$\pm$.071} & \underline{.679}\,{\scriptsize$\pm$.129} & \underline{.721}\,{\scriptsize$\pm$.041} & \underline{.854}\,{\scriptsize$\pm$.021} & \underline{.714}\,{\scriptsize$\pm$.024} & .746\,{\scriptsize$\pm$.013} & \textbf{.751}\,{\scriptsize$\pm$.019} & \underline{.729}\,{\scriptsize$\pm$.017} & \textbf{.778}\,{\scriptsize$\pm$.013} \\
 & Value probe & \textbf{.717}\,{\scriptsize$\pm$.031} & \textbf{.692}\,{\scriptsize$\pm$.030} & \textbf{.737}\,{\scriptsize$\pm$.004} & .814\,{\scriptsize$\pm$.010} & \textbf{.731}\,{\scriptsize$\pm$.003} & \textbf{.761}\,{\scriptsize$\pm$.010} & \underline{.744}\,{\scriptsize$\pm$.011} & \textbf{.733}\,{\scriptsize$\pm$.019} & .712\,{\scriptsize$\pm$.026} \\
\midrule
\multirow{4}{*}{gpt-oss-20b} & Confidence & .542\,{\scriptsize$\pm$.023} & .529\,{\scriptsize$\pm$.016} & .549\,{\scriptsize$\pm$.050} & .528\,{\scriptsize$\pm$.016} & .541\,{\scriptsize$\pm$.023} & .549\,{\scriptsize$\pm$.050} & .555\,{\scriptsize$\pm$.045} & .525\,{\scriptsize$\pm$.009} & .525\,{\scriptsize$\pm$.016} \\
 & Probe & \underline{.812}\,{\scriptsize$\pm$.021} & \textbf{.692}\,{\scriptsize$\pm$.016} & .677\,{\scriptsize$\pm$.018} & \underline{.797}\,{\scriptsize$\pm$.011} & \underline{.651}\,{\scriptsize$\pm$.045} & .596\,{\scriptsize$\pm$.021} & .713\,{\scriptsize$\pm$.033} & \underline{.739}\,{\scriptsize$\pm$.055} & \underline{.651}\,{\scriptsize$\pm$.018} \\
\cmidrule(lr){3-5}\cmidrule(lr){6-8}\cmidrule(lr){9-11}
 & Counterf.\ sim. & \textbf{.828}\,{\scriptsize$\pm$.099} & \underline{.652}\,{\scriptsize$\pm$.144} & \textbf{.759}\,{\scriptsize$\pm$.034} & \textbf{.799}\,{\scriptsize$\pm$.088} & \textbf{.670}\,{\scriptsize$\pm$.168} & \textbf{.624}\,{\scriptsize$\pm$.057} & \textbf{.796}\,{\scriptsize$\pm$.027} & \textbf{.784}\,{\scriptsize$\pm$.092} & \textbf{.660}\,{\scriptsize$\pm$.091} \\
 & Value probe & .743\,{\scriptsize$\pm$.010} & .636\,{\scriptsize$\pm$.028} & .645\,{\scriptsize$\pm$.028} & .700\,{\scriptsize$\pm$.007} & .633\,{\scriptsize$\pm$.023} & .591\,{\scriptsize$\pm$.007} & .697\,{\scriptsize$\pm$.005} & .721\,{\scriptsize$\pm$.018} & \underline{.653}\,{\scriptsize$\pm$.015} \\
\bottomrule
\end{tabular*}
\end{table*}


% Numbers (2026-10-27): generated/tab_verdict_summary.tex,
% tab_language_detail_verdict.tex, paired_shift.csv (fold=mean, train=all: mean
% over the six language pairs, three partitions; "±" is 1.96 se_total). Per-pair
% statements are descriptive (fold 0/1/2 rows of paired_shift.csv, value probe
% − noncausal): ahead in all three partitions on pyctl→java, cpp→pyctl,
% cpp→java, java→pyctl; behind in all three on java→cpp. Within-problem AUROC
% from layer_selection{,_f1,_f2}.csv (mean over languages, then partitions);
% figure count from the score layer (artifact/manifest.csv): value-probe
% checkpoints above the probe's best layer on the shifted rows, Qwen3-Coder
% L35--L39 in each fold.

%The verdict target asks whether the model's judgment of a candidate can be
%trusted (Section~\ref{sec:problem}): the decision an agentic pipeline faces
%when it can either accept the verdict or run the tests. A high AUROC means that
%a method ranks candidates with a correct verdict above those with a wrong one.
Both LLMs are mediocre judges, with balanced accuracies of their verdicts of 57--74\% (Table~\ref{tab:datasets}) against the tests. RQ1 asks how well our methods identify the candidates on which the model's verdict is correct.
Table~\ref{tab:verdict-summary} summarizes the results, showing the mean per method over partitions.
Table~\ref{tab:language-detail-verdict} breaks down the language shift by training and evaluation language.

\paragraph{In distribution, hidden states outperform confidence}
On LCB and the language sets, the methods that read hidden states reach 0.69--0.83, while confidence is weak for
Qwen3-Coder (0.64--0.67) and close to chance for gpt-oss (0.54--0.57). %, as reported for natural-language tasks~\cite{Kadavath2022LanguageM, OrgadTGRSKB25}. 
On the language sets, the probe and the two causal methods are
level for Qwen3-Coder (0.758--0.761); for gpt-oss, counterfactual simulation is best (0.808). BCB is harder, no trained method exceeds 0.639 for Qwen3-Coder or 0.601 for gpt-oss.

\paragraph{Under language shift, the causal subspace is the most robust predictor for Qwen3-Coder}
%Probes on hidden states can degrade under distribution shift~\cite{OrgadTGRSKB25, Kramar2026Probes}. 
For Qwen3-Coder, the probe drops from 0.76 in distribution to 0.67 when the language changes%(mean over the six pairs)
, while both causal methods fall only from 0.76 to 0.73, against 0.64 for confidence. Averaged over the six language pairs, the value probe is ahead of
the best non-causal method of each pair by $0.042 \pm 0.021$ (95\% interval of the paired test, Section~\ref{sec:setup-protocol}) and ahead in each of the three partitions ($+0.035$, $+0.053$, $+0.038$). Counterfactual simulation
reaches the same mean but varies more between partitions ($\pm 0.043$). 
%Per pair, the value probe is ahead in all three partitions on four of the six pairs and behind in all three on Java$\to$C++, C++ being the hardest target language. 
The lead does not depend on the probe's layer choice: in 13 of the 18 pair--partition combinations, most value-probe checkpoints lie above the probe at its best layer on the shifted rows, as illustrated for Qwen3-Coder, trained on C++ in Figure~\ref{fig:verdict-profile}. 
%These results suggest that the learned subspace provides a more robust signal of verdict correctness than the full hidden state for Qwen3-Coder. 
For gpt-oss, all trained methods lose 0.07--0.12 AUROC under language shift, and neither causal method is significantly ahead of the best non-causal one. %Counterfactual simulation is ahead in five of the six pairs, but its lead varies between partitions from $-0.06$ to $+0.11$. 
The robustness reported for natural language tasks~\cite{Huang2025Internal} thus carries over to code for one of the two models; Section~\ref{sec:rq-subspace} examines what carries the gain.

% Figure: scripts/paper/fig_layer_profile.py --target verdict --cells qwen3:cpp
%   --wide --first-block 30 --out paper/generated/fig_verdict_profile_qwen3_cpp.pdf
% (fold 0, score layer; one row since 2026-10-28, probe below block 30 at
% 0.49--0.64 not drawn).
\begin{figure}[t]
\centering
\includegraphics[width=\columnwidth]{generated/fig_verdict_profile_qwen3_cpp}
\caption{AUROC per layer on the verdict target, Qwen3-Coder trained on C++
(fold 0), in distribution (left) and under language shift. Rings mark the
layer chosen in distribution; the dashed line is confidence. Under shift to
Python, every causal checkpoint lies above the probe's best layer; on Java,
most do.}
\label{fig:verdict-profile}
\end{figure}

% Figure: scripts/paper/make_result_tables.py --only fw (layer_selection{,_f1,_f2}.csv,
% columns pooled and within; 2026-10-27).
\begin{figure}[t]
\centering
\includegraphics[width=\columnwidth]{generated/fig_within}
\caption{Within-problem (hollow) against pooled AUROC (filled), in distribution
(faint) and under language shift: means over training languages and partitions.
Within-problem AUROC uses the problems with mixed outcomes only
(Section~\ref{sec:setup-protocol}); the length of a line is how much of the
pooled AUROC comes from ranking problems.}
\label{fig:within}
\end{figure}

\paragraph{The advantage lies in identifying unreliable problems, not candidates}
Pooled AUROC combines the ability to identify problems on which the judge is unreliable with the ability to distinguish correct from incorrect verdicts within a problem (Section~\ref{sec:setup-protocol}). The second is limited. Under language shift, within-problem AUROC is 0.616--0.618 for all trained methods of Qwen3-Coder. 
For gpt-oss the probe reaches 0.652, counterfactual simulation 0.635 and the value probe 0.576 (Figure~\ref{fig:within}).
The causal methods' pooled advantage thus comes mainly from identifying problems on which the judge is unreliable, rather than individual candidates it misjudges. 

\paragraph{A change of benchmark affects every method more than a change of
language}
Transfer between LCB %, whose problems are algorithmic, 
and BCB %, whose tasks rely on libraries, 
is asymmetric (Table~\ref{tab:verdict-summary}). From LCB to BCB, all trained methods of Qwen3-Coder fall
to 0.453--0.485, below confidence (0.552). In the reverse direction, the probe stays near chance (0.516), while the causal methods reach 0.680 and 0.695, close to confidence (0.669).
For gpt-oss, every method stays at 0.48--0.61 in both directions. %With one partition, these differences between methods are descriptive. 
%\orange{Section~\ref{sec:rq-subspace} relates the failure to the model's verdicts on BCB. }

\begin{rqanswer}{RQ1}
In distribution, hidden states predict whether the model's verdict is right far
better than its confidence. Under language shift, a probe on the causally
localized subspace is the most robust predictor for Qwen3-Coder (+0.04 AUROC over
the best non-causal method, in all three partitions), mainly by identifying problems
on which the judge is unreliable; for gpt-oss, no method is significantly ahead. 
Verdict prediction does not transfer reliably between benchmarks.
\end{rqanswer}


\subsection{RQ2: Which part of the causal methods carries their signal?}
\label{sec:rq-subspace}

On Qwen3-Coder's verdict target, the causal methods are the most robust predictors under language shift (RQ1). 
We now investigate which aspects of the method contribute to this performance: whether prediction requires the intervention itself, whether the learned directions matter beyond their dimensionality, and whether interchange intervention accuracy (see Section~\ref{sec:methods-causal}) indicates when the subspace is useful.


% Figure: scripts/paper/fig_subspace_layers.py (fold 0, results/*_randomlayers
% and the block-fixed learned value probes; 2026-10-22). Space permitting; it
% could replace Table~\ref{tab:subspace-control}.
\begin{figure}[t]
\centering
\includegraphics[width=\columnwidth]{generated/fig_subspace_layers}
\caption{Value probe on the learned DAS subspace against value probes on
three random 6-dimensional subspaces at the same layer, verdict target,
mean over shifted languages (the band is the range over
the random subspaces). Shaded: layers where DAS has
converged (IIA $\geq$ 0.9). On Qwen3-Coder the learned subspace is above
random at every layer; on gpt-oss it is above random before convergence and
near chance once DAS has converged.}
\label{fig:subspace-layers}
\end{figure}

\paragraph{Reading the subspace is sufficient}
Counterfactual simulation intervenes through the subspace and reads the answer
token, at the cost of $B$ additional forward passes per candidate; the value
probe reads the subspace directly in the judging pass. On Qwen3-Coder the two are equally robust under language shift (mean difference over the six pairs $-0.000 \pm 0.033$ on the verdict target, $+0.004 \pm 0.007$ on the program target), and both are ahead of confidence, which reads the same
answer token without the subspace ($+0.085 \pm 0.033$ and $+0.05 \pm 0.03$).
%Preliminary experiments showed that the choice of the $B=16$ bases barely matters: drawing them from different problems changed the pooled AUROC of one checkpoint by at most 0.01. 
The intervention thus brings no significant improvement over reading the
subspace, and the value probe costs no more at prediction time than a probe.

\paragraph{The learned direction matters, not the low dimension}
A six-dimensional readout might transfer better simply because it is small. We
therefore train the value probe on three random subspaces of the same
dimension at the same layer. At the layer chosen in distribution, the
learned subspace reaches 0.70--0.75 under shift for Qwen3-Coder, random
subspaces 0.48--0.61, no better than the probe on the full hidden state (0.668,
Table~\ref{tab:verdict-summary}). The same holds at every trained layer (Figure~\ref{fig:subspace-layers}). 
%\orange{Die Figure (insb. untere Zeile) wirft mehr Fragen auf... Auf Qwen3-Coder fokussierene evtl.?}

% % Undecided (2026-10-27): why the subspace helps on the verdict target;
% % hypotheses and checks in PAPER_PLAN_FORGE.md, open decision 14.
% \paragraph{\orange{The causal score reads how firmly a verdict is held}}
% \orange{The verdict score of counterfactual simulation is the probability that
% the bases output the candidate's verdict once its subspace is patched in. A
% base the patch does not move keeps its own answer, so the score carries an
% offset between \texttt{Yes} and \texttt{No} verdicts that depends on the bases,
% not on the candidate. Within one verdict the offset is constant, and there
% right verdicts score higher than wrong ones in every setting: under language
% shift, the AUROC among accepted and among rejected candidates is 0.65 and 0.77
% for Qwen3-Coder and 0.64 and 0.68 for gpt-oss (means over pairs and
% partitions). The causal methods thus read how firmly a candidate's state holds
% the model's decision. The offset explains the failure under benchmark shift:
% from LCB to BCB, Qwen3-Coder's \texttt{Yes} verdicts score 0.77 higher than its
% \texttt{No} verdicts, and 84\% of its wrong verdicts there are \texttt{Yes}, so
% the pooled AUROC falls to 0.43, although right verdicts still score higher within
% each verdict (0.57 and 0.61).}

\paragraph{Interchange accuracy does not guarantee predictive utility}
DAS learns the subspace so that swapping it between two inputs swaps the
model's verdict. IIA measures how often this succeeds. A high IIA thus shows
that the subspace controls the decision, not that it tells whether the decision
is right. On gpt-oss the two diverge (Figure~\ref{fig:subspace-layers}): at L11
and L13, where DAS succeeds only partly (IIA 0.54--0.64), the learned subspace
transfers better than random ones (0.56--0.68 against 0.52--0.55 under shift),
and L13 is the layer chosen in distribution. However, at L17--L21, where DAS almost
always succeeds (IIA 0.91--0.95), it falls to 0.52--0.58, no better than
random. %\orange{One reading is that the decision is fully formed at these layers, so the subspace holds a clean \texttt{Yes} or \texttt{No} that says little about whether it is right.} 
On Qwen3-Coder the converged layers (L39--L43, IIA 0.95) remain informative. IIA thus cannot tell the two cases apart, and layers are chosen by held-out prediction instead (Section~\ref{sec:setup-training}).

\begin{rqanswer}{RQ2}
Where the causal methods help, the learned subspace carries the gain: the value
probe, which reads it in the judging pass, is as robust as counterfactual
simulation, and random subspaces of the same size transfer only as well as the
full hidden state. %\orange{Within each verdict, right verdicts are held more firmly in the subspace than wrong ones; across verdicts, counterfactual simulation adds an offset set by its bases, which explains its failure from LCB to BCB.} 
IIA does not indicate when the subspace is useful: on gpt-oss, the
converged subspaces are no better than random ones.
\end{rqanswer}
 

\subsection{RQ3: How well do the methods predict whether the program is correct?}
\label{sec:rq-program}


The program target asks whether a candidate passes its tests, regardless of
the model's verdict. The two targets
agree only where the judge is right. %Across all candidates, they are unrelated: the model judges passing and failing programs equally accurately. %(Adrian: die zwei Sätze mag ich nicht. Sicher dass korrekt?)}
We therefore ask whether the results on the verdict target carry over. Table~\ref{tab:program-summary} summarizes the results. In distribution, all methods predict
program correctness well above chance, AUROC ranges from 0.70--0.88 on LCB and the language
sets, with the hidden-state probe performing best on the language sets for both models. BCB is more challenging for every method.

% Generated by scripts/paper/make_result_tables.py -- do not edit.
\begin{table}[t]
\centering\footnotesize
\setlength{\tabcolsep}{3.2pt}
\caption{AUROC for program correctness (mean over 3 partitions). \emph{LCB}, \emph{BCB}: trained and evaluated on that benchmark; $\to$: evaluated on the other one. Language shift: \emph{ID} is the mean over the three training languages, \emph{Shift} the mean over the six train$\to$evaluation pairs, $\Delta$ their difference. %In-distribution cells are cross-fitted over problems; shifted cells use the layer chosen in distribution. 
Best per column and model in bold.}
\label{tab:program-summary}
\begin{tabular}{l rrrr rrr}
\toprule
 & \multicolumn{4}{c}{Benchmark shift} & \multicolumn{3}{c}{Language shift} \\
\cmidrule(lr){2-5}\cmidrule(lr){6-8}
Method & LCB & $\to$BCB & BCB & $\to$LCB & ID & Shift & $\Delta$ \\
\midrule
\multicolumn{8}{l}{\emph{Qwen3-Coder-30B}} \\
Confidence & .815 & .650 & .648 & .815 & .759 & .755 & \textbf{$-$.003} \\
Probe & \textbf{.844} & \textbf{.656} & \textbf{.688} & .697 & \textbf{.878} & .790 & $-$.088 \\
LAT & .844 & .580 & .590 & .731 & .807 & .799 & $-$.009 \\
\addlinespace[2pt]
Counterf.\ sim. & .839 & .644 & .635 & .813 & .848 & .806 & $-$.042 \\
Value probe & .838 & .642 & .664 & \textbf{.815} & .861 & \textbf{.810} & $-$.051 \\
\midrule
\multicolumn{8}{l}{\emph{gpt-oss-20b}} \\
Confidence & .723 & \textbf{.648} & .646 & .721 & .696 & .691 & \textbf{$-$.004} \\
Probe & .793 & .575 & .649 & .690 & \textbf{.859} & .693 & $-$.166 \\
LAT & .767 & .618 & .608 & .676 & .760 & .666 & $-$.094 \\
\addlinespace[2pt]
Counterf.\ sim. & .780 & .647 & .652 & \textbf{.724} & .820 & \textbf{.710} & $-$.111 \\
Value probe & \textbf{.803} & .638 & \textbf{.652} & .712 & .846 & .666 & $-$.180 \\
\bottomrule
\end{tabular}
\end{table}



\paragraph{Under language shift, the causal methods lose their advantage}
The full-state probe loses 0.088 AUROC for Qwen3-Coder and 0.166 for gpt-oss, while confidence changes little ($-0.003$ and $-0.004$). For Qwen3-Coder, LAT and both causal methods nevertheless remain significantly ahead of confidence by 0.044--0.054 AUROC on average; the probe's lead of 0.035 is no longer significant. For gpt-oss, no trained method significantly outperforms confidence. Importantly, neither causal method outperforms the best non-causal method: for Qwen3-Coder, the value probe and counterfactual simulation trail it by 0.010 and 0.014 AUROC, respectively, neither significantly. For gpt-oss, the value probe is significantly behind ($-0.068$), while counterfactual simulation is statistically level ($-0.025$). Thus, the robustness of causal localization observed for verdict correctness (RQ1) does not carry over to program correctness.


% In distribution, the trained methods are 0.05--0.16 ahead of confidence on the
% language sets. Under language shift the probe loses most ($-0.088$ for
% Qwen3-Coder, $-0.166$ for gpt-oss), while confidence barely changes ($-0.003$,
% $-0.004$). For Qwen3-Coder, LAT and the two causal methods remain
% significantly ahead of confidence, by 0.044--0.054 on average over the pairs;
% the probe's lead ($+0.035 \pm 0.035$) is no longer significant. For gpt-oss,
% no trained method is ahead ($-0.037$ to $+0.019$, none significant). Where no labeled data exists for the new language, confidence is thus about 0.05 behind the best
% trained method for Qwen3-Coder and as good as any for gpt-oss.

% \paragraph{The causal methods have no advantage on program correctness}
% On average over the pairs, neither causal method is ahead of the best
% non-causal method under language shift: for Qwen3-Coder, the differences are
% $-0.010 \pm 0.019$ (value probe) and $-0.014 \pm 0.017$ (counterfactual
% simulation), neither significant; for gpt-oss, the value probe is
% significantly behind ($-0.076 \pm 0.035$) and counterfactual simulation is
% level ($-0.025 \pm 0.034$). Their robustness on the verdict target (RQ1) does
% not carry over to program correctness.

\paragraph{Program correctness transfers between benchmarks}
In contrast to verdict correctness, predictors for program correctness transfer between LCB and BCB in both directions. When trained on LCB and evaluated on BCB, methods reach 0.580--0.656 for Qwen3-Coder, and 0.575--0.648 for gpt-oss, comparable to
the respective in-distribution results on BCB (0.590--0.688 and 0.608--0.652).
 Transfer in the reverse direction is similarly strong, reaching 0.70--0.82 for Qwen3-Coder and 0.68--0.72 for gpt-oss.
% Figure removed for space (2026-10-23): the selection cost is given in the
% text. scripts/paper/fig_layer_profile.py --target program --cells qwen3:cpp
% gptoss:cpp
% \begin{figure}[t]
% \centering
% \includegraphics[width=\columnwidth]{generated/fig_program_profile_qwen3_cpp_gptoss_cpp}
% \caption{AUROC per layer on the program target, in distribution (top) and under
% language shift (fold 0), trained on C++: Qwen3-Coder (left), gpt-oss
% (right). Rings mark the layer chosen in distribution; the dashed line is
% confidence at its chosen temperature. The layer that is best for the probe in
% distribution is not the one that transfers. LAT is omitted: its per-layer sign
% fit flips in many layers.}
% \label{fig:layers}
% \end{figure}
\paragraph{Pooled AUROC overstates candidate-level discrimination}
As for verdict correctness, pooled AUROC reflects both the ability to identify problems the model tends to solve and the ability to distinguish candidates within a problem.
Choosing a candidate needs only the second, and here every method is weak (Figure~\ref{fig:within}, bottom): under language shift, within-problem AUROC lies at 0.58--0.64 for all methods and both models, against a pooled AUROC of 0.67--0.81. For Qwen3-Coder confidence performs best within problem (0.629). 
The trained methods reach higher within-problem AUROCs of 0.69--0.78 only for gpt-oss in distribution, but lose this advantage under shift. Consequently, the predictors are more useful for identifying problems that warrant verification than for choosing among candidates of the same problem.
Section~\ref{sec:rq-problem} predicts at problem level directly.

% Numbers below: generated/tab_genprobe.tex from generated/genprobe_compare.csv
% (scripts/analysis/genprobe_compare.py, fold 0, 2026-10-22).
\paragraph{Reading the generation is as good as asking}
All methods so far read the model in a separate judging pass. The same information might already be present when the model generates the candidate. %\orange{Verstehe ich nicht. Ist das nicht per Definition identisch?}. 
We therefore train the same probe on the hidden states of the model's own
generation (without the judgment prompt), replayed token by token, and compare it with the judging-pass probe on identical candidates, Table~\ref{tab:genprobe} reports the results. We compare different token positions here: at the last token of the candidate's code, it is
at least as good as the judging pass in distribution (0.895 against 0.886 for
Qwen3-Coder, 0.864 against 0.851 for gpt-oss); under language shift, it is
close for Qwen3-Coder (0.795 against 0.779) and behind for gpt-oss (0.673
against 0.704). A correctness score can thus be read while the code is
written. Read at the last prompt token, before any code exists, the probe
has seen only the problem, yet for Qwen3-Coder it matches the judging pass in
distribution (0.891). Much of what the methods predict is thus the problem,
not the candidate.

% Generated by scripts/paper/make_result_tables.py -- do not edit.
\begin{table}[t]
\centering\footnotesize
\setlength{\tabcolsep}{3.2pt}
\caption{The probe on the judging pass against the same probe on the states of the model's own generation, program target on the language sets. \emph{ID}: mean AUROC over the three languages, \emph{Shift}: mean over the six language pairs, \emph{Within}: AUROC over pairs of candidates of the same problem. The last prompt token has seen only the problem.}
\label{tab:genprobe}
\begin{tabular}{l rr rr}
\toprule
 & \multicolumn{2}{c}{ID} & \multicolumn{2}{c}{Lang. Shift} \\
\cmidrule(lr){2-3}\cmidrule(lr){4-5}
Readout & Pooled & Within & Pooled & Within \\
\midrule
\multicolumn{5}{l}{\emph{Qwen3-Coder-30B}} \\
Judging pass & .886 & .553 & .779 & .540 \\
Generation, last code token & .895 & .620 & .795 & .627 \\
Generation, last prompt token & .891 & .514 & .755 & .531 \\
\midrule
\multicolumn{5}{l}{\emph{gpt-oss-20b}} \\
Judging pass & .851 & .772 & .704 & .609 \\
Generation, last code token & .864 & .781 & .673 & .623 \\
Generation, last prompt token & .720 & .524 & .669 & .478 \\
\bottomrule
\end{tabular}
\end{table}


\begin{rqanswer}{RQ3}
All methods predict program correctness well above chance, largely by recognizing the problems the model fails on, %Within one problem, every method is weak 
and this transfers better between benchmarks than the verdict. Under
language shift, the trained methods lose most of their lead over confidence:
about 0.05 remains for Qwen3-Coder, nothing for gpt-oss, and the causal
methods' robustness on the verdict target does not carry over. A probe on the model's own generation is as good as one on a separate judging pass in distribution.
\end{rqanswer}


\subsection{RQ4: Does agreement across samples improve problem-level prediction?} % tell better than a single pass which problems the model will solve?}
\label{sec:rq-problem}
% Numbers from generated/tab_problem_level.tex and generated/problem_level_ci.csv
% (scripts/analysis/problem_level_ci.py, 2026-10-16). "Significant": 95% paired
% problem bootstrap (2000 draws, same rows and averaging as the table) excludes
% zero. Dropped 2026-10-28: the before-sampling paragraph (problem-only probe
% now in RQ3), the end-of-reasoning probe, per-pair counts, the Harrell's C
% note, the sentence on prior consistency studies (covered in related work).

This question moves from the candidate to the problem level. Much of the
candidate-level signal reflects which problems the model solves
(Section~\ref{sec:rq-program}). When several candidates are sampled for a
problem, their agreement provides a further signal.
EigenScore~\cite{Chen24Eigenscore} reads this agreement from the hidden states
of the generations (Section~\ref{sec:methods-consistency}). Like confidence, it
needs no training, so it cannot lose under shift. We compare at equal sampling
cost: each single-pass method is averaged over the candidates EigenScore uses,
and every method predicts whether at least half of a problem's candidates pass
(Table~\ref{tab:problem-level}).

% - A possibility to prevent the problems under distribution shift, is to use an unsupervised method, that needs no training. Also other advantages, no labeled training data needed\dots
% One such method is eigenscore, which uses cnsistency across the hidden states of n different samples (n=5 in our experiments) to compute a score for 'how likely that this problem will be solved by the model'. As a consistency-based method, eigenscore scores at problem level in contrast to our other candidate level methods. To compare it to our methods nonetheless, we mean-pool our candidate scores to problem scores and label them by majority vote -> a problem is solved if more than $50\%$ of its candidates pass (reasonable since most (X-Y\% depending on the benchmark) of our problems contain passing or failing candidates only).

% Remember that we exclude empty generations (where no program could be extracted from the models reasoning) as trivial from our experiments, so we filter out all problems with less than n=5 non-empty samples, resulting in a slightly reduced test set for this experiment.

% Generated by scripts/paper/make_result_tables.py -- do not edit.
\begin{table}[t]
\centering\footnotesize
\setlength{\tabcolsep}{3.2pt}
\caption{AUROC for whether a problem is solved (at least half of its candidates pass). Single-candidate methods are averaged over the candidates of a problem, at the layer of Table~\ref{tab:program-summary}; EigenScore is computed from the sampled candidates and needs no training. Columns as in Table~\ref{tab:program-summary}. $\uparrow$: significantly above EigenScore (paired problem bootstrap, 95\%; for ID and Shift on the mean over languages). All methods on the problems that have an EigenScore.$^{*}$}
\label{tab:problem-level}
\begin{tabular}{l rrrr rr}
\toprule
 & \multicolumn{4}{c}{Benchmark shift} & \multicolumn{2}{c}{Language shift} \\
\cmidrule(lr){2-5}\cmidrule(lr){6-7}
Method & LCB & $\to$BCB & BCB & $\to$LCB & ID & Shift \\
\midrule
\multicolumn{7}{l}{\emph{Qwen3-Coder-30B}} \\
EigenScore & .504 & .658 & .658 & .504 & .528 & .528 \\
Confidence & .877$^{\uparrow}$ & \textbf{.670} & .671 & .872$^{\uparrow}$ & .774$^{\uparrow}$ & .775$^{\uparrow}$ \\
Probe & .906$^{\uparrow}$ & .638 & .671 & .884$^{\uparrow}$ & \textbf{.911}$^{\uparrow}$ & .803$^{\uparrow}$ \\
LAT & .902$^{\uparrow}$ & .618 & \textbf{.682} & .871$^{\uparrow}$ & .844$^{\uparrow}$ & .821$^{\uparrow}$ \\
\addlinespace[2pt]
Counterf.\ sim. & \textbf{.912}$^{\uparrow}$ & .646 & .676 & .881$^{\uparrow}$ & .875$^{\uparrow}$ & .831$^{\uparrow}$ \\
Value probe & .904$^{\uparrow}$ & .652 & .680 & \textbf{.889}$^{\uparrow}$ & .890$^{\uparrow}$ & \textbf{.838}$^{\uparrow}$ \\
\midrule
\multicolumn{7}{l}{\emph{gpt-oss-20b}} \\
EigenScore & .660 & .518 & .518 & .660 & .586 & .586 \\
Confidence & .698 & \textbf{.677}$^{\uparrow}$ & .678$^{\uparrow}$ & .692 & .623 & .615 \\
Probe & .791 & .615 & \textbf{.712}$^{\uparrow}$ & \textbf{.759} & \textbf{.845}$^{\uparrow}$ & \textbf{.676}$^{\uparrow}$ \\
LAT & \textbf{.808}$^{\uparrow}$ & .662$^{\uparrow}$ & .622$^{\uparrow}$ & .688 & .772$^{\uparrow}$ & .645 \\
\addlinespace[2pt]
Counterf.\ sim. & .772 & .655$^{\uparrow}$ & .685$^{\uparrow}$ & .692 & .800$^{\uparrow}$ & .662 \\
Value probe & .790 & .663$^{\uparrow}$ & .692$^{\uparrow}$ & .676 & .834$^{\uparrow}$ & .630 \\
\bottomrule
\end{tabular}
\par\smallskip
\parbox{\linewidth}{\scriptsize $^{*}$EigenScore needs $m$ non-empty samples: five per problem on LCB and the language sets, ten on BCB. Problems with fewer non-empty generations are left out for every method, and the single-candidate methods are averaged over the candidates of the same samples. This keeps 93--100\% of the evaluation problems for Qwen3-Coder-30B and 76--94\% for gpt-oss-20b. %Of the problems left out (counted once per benchmark), 0 of 40 are solved for Qwen3-Coder-30B and 47 of 143 for gpt-oss-20b.
}
\end{table}


% Numbers: generated/tab_genprobe.tex (candidate-level pooled AUROC, fold 0).
% A problem-level row for the last-prompt-token probe in
% Table~\ref{tab:problem-level} would make this directly comparable with
% EigenScore (make_result_tables.py t5); not done (2026-10-28).
The cheapest prediction needs no sample at all: for Qwen3-Coder, a probe that
has seen only the problem %\orange{verstehe ich nicht} 
already ranks candidates as well as the judging pass
(Section~\ref{sec:rq-program}). Sampling several candidates pays off for
prediction only if their agreement adds to what a single pass provides.

\paragraph{Agreement provides no consistent advantage}
For Qwen3-Coder, EigenScore is at chance on LCB (0.504) and on the language
sets (0.528), where every single-candidate method, confidence included, is
significantly ahead (0.77--0.91). On BCB it is comparable to them (0.658 against
0.671--0.682). For gpt-oss the benchmarks swap roles: EigenScore is at chance
on BCB (0.518 against 0.622--0.712) and closer on LCB (0.660 against
0.698--0.808). Under shift, the single-pass methods lose ground. 
From LCB to BCB, Qwen3-Coder's trained methods fall slightly below EigenScore (0.618--0.652 against 0.658) while confidence stays above it (0.670); under language shift, the
gap remains large for Qwen3-Coder (0.775--0.838 against 0.528) and narrows for
gpt-oss (0.615--0.676 against 0.586).


\begin{rqanswer}{RQ4}
Agreement across samples gives no consistent advantage over a single pass.
EigenScore is level with the single-pass methods in some cells and near
chance in others, and confidence, which needs no labels either, is never
significantly below it. That EigenScore cannot lose under shift does not help where it starts near chance.
\end{rqanswer}

% Notes from earlier drafts:
% RQ2: How well can they tell whether the model's own verdict on its code can
% be trusted, and how much of that survives shift?
% How does an unsupervised consistency-based method compare to these methods?
% Numbers are reported in mean±std. For each task, the highest AUROC value is
% in bold. Values that are within a standard deviation are underlined.


\section{Discussion}
\label{sec:discussion}
%- Verdict, In distribution: The practical takeaway here is that for monitoring a judge, reading its internal states is more reliable than asking it or reading its confidence.
% Verdict shift: They are consequently more useful for prioritizing problems for verification than for selecting among candidates of the same problem.
%triage instead of selection since they help prioritize which problems to verify rather than select among a problem's candidates.

\paragraph{Verdict and program are different prediction targets}
Our results show that the usefulness of execution-free correctness prediction depends on both the prediction target and the type of distribution shift. In particular, predicting whether a model's verdict is correct is different from predicting whether the program itself is correct, even though both targets use signals from the same judging pass.
The two targets respond differently to shift. Causal localization helps only on the verdict, and between benchmarks only program correctness transfers. An execution-free correctness predictor for an LLM-based code-generation system should therefore be chosen and evaluated for the decision it supports. To decide which candidates to test in a new programming language, output confidence is a strong baseline: it needs no labels and loses almost nothing under shift,
while trained predictors keep only part of their lead. To decide whether to
trust the model's verdict, confidence is weak, and no verdict predictor we tested
carries over to a new benchmark. % without data from it.


% \paragraph{What a practitioner can take away}
% To decide whether an LLM's verdict on code can be trusted in a new
% programming language, a probe on a causally localized subspace was the most
% robust predictor we found, and once the subspace is learned it costs no more
% than any probe. The gain is model-dependent, however, and DAS's own quality
% measure does not reveal it: on gpt-oss, well-converged subspaces were no
% better than random ones. A practitioner therefore has to validate the subspace
% on held-out data from the target setting rather than trust IIA. Between
% benchmarks, no verdict predictor transferred reliably. The judge's errors
% change character: Qwen3-Coder accepts 58\% of the LCB candidates, of which
% 55\% pass, but 72\% of the BCB candidates, of which only 40\% pass
% (Table~\ref{tab:datasets}), so a predictor of its reliability learned on one
% benchmark meets a different error profile on the other. A few labeled
% examples from the target setting, which help probes in other domains
% considerably~\cite{McKenzie2025HighStakes}, are likely needed there.
% For program correctness, output confidence is a strong baseline that needs no
% labels and loses almost nothing under language shift; trained probes pay off
% in distribution, and much less after a shift.

\paragraph{Triage rather than selection}
Much of what the methods predict is which problems
the model solves: every method is weak at ranking candidates within a problem, and a probe that has seen only the problem matches one that has seen the candidate (Table \ref{tab:genprobe}). Pooled metrics,  such as AUROC~\cite{bouchard2026functional, shi2026codemorethantext} or accuracy and F1~\cite{bui2025openia, Vu2026Autoprobe}, thus mix problem difficulty with the ability to tell a correct candidate from an incorrect one. 
For triage across problems this is useful, while for choosing among
the candidates of one problem it is not. Selection metrics~\cite{ribeiro2026repe} and pairwise
accuracy within a task~\cite{ribeiro2026robustness} measure the second ability
directly; we recommend reporting such a within-problem measure alongside the
pooled one. 

\paragraph{Causal localization for code}
Huang et al.~\cite{Huang2025Internal} report that causally localized
predictors are more robust under shift on natural-language tasks. For code,
this holds on one of our two models and for the target the method was
designed for, the model's own verdict; on the program correctness target it gives no
advantage. Where it helps, the learned subspace carries the gain: once the subspace is
learned, a probe on it is as robust as simulating interventions and costs no
more than any probe. IIA does not reliably identify subspaces that are useful for correctness prediction. On gpt-oss, well-converged subspaces were no better than random ones, so the
subspace has to be validated on held-out data.
%, not the intervention: reading the subspace with a probe is as robust as simulating interventions, at a fraction of the cost. %\orange{Why the subspace helps on Qwen3-Coder and not on gpt-oss we can only partly explain. On gpt-oss, the useful subspaces are those at layers where DAS has not converged; one possibility is that gpt-oss, a reasoning model judging without reasoning, has fully formed its decision at these layers, so that the subspace holds the decision but little about whether it is right.}

\section{Threats to Validity}
\label{sec:threats}
% Restructured 2026-10-29 (construct / internal / external). Reasoning numbers:
% same 498 samples for both settings, none on LCB-Java; low effort (median 150
% tokens) gives balanced accuracy 63.9% (not reported). Removed: "each
% comparison uses one machine" (rule dropped) and the prose-program-rows
% sensitivity sentence (Viola, 2026-10-29). Python header: FINDINGS
% 2026-09-22, effect not measured.
\paragraph{Construct validity}
Correctness is defined by the benchmarks' tests, which can miss errors. The
judging prompt asks for a single token without reasoning; gpt-oss, a
reasoning model, judges moderately  better with reasoning: on 498 program-target
candidates of the language sets, medium-effort reasoning (median 1{,}101
tokens) raises its balanced accuracy from 62.4\% to 67.8\% and the AUROC of its
\texttt{Yes} probability from 0.68 to 0.73, still below the trained internal
methods (Table~\ref{tab:program-summary}), at the cost of the reasoning
tokens. %The judging prompt introduces every task as a Python task, also for the C++ and Java sets; this may lower the judge's accuracy there, but all methods except LAT read the same judging pass, so the comparison between methods is unaffected. LAT reads its published stimulus prompt (Section~\ref{sec:methods}). 
%EigenScore is computed on non-empty generations only; keeping the (trivial) empty generations raises it to 0.597 on LCB and 0.610--0.646 on the language sets for Qwen3-Coder and to 0.719 on LCB for gpt-oss, still below the best single-pass method in each of these cells.

\paragraph{Internal validity}
The trained methods choose among different numbers of layers: the probe and
LAT among all 25--49, the causal methods among the 4--7 layers at which DAS
was trained, since each subspace takes several GPU hours. 
A wider choice can find a better layer, but under shift the layer best in distribution need not transfer, so the effect can go either way. Under shift, however, the causal checkpoints often lie above the probe even at the probe's best layer (Figure~\ref{fig:verdict-profile}).
Counterfactual simulation uses one fixed set of 16 bases;
drawing them from different problems changed the AUROC of one checkpoint by at
most 0.01 AUROC. 
Reruns of the same configuration on different machines differ by up to 0.02 AUROC on identical rows, so we do not interpret smaller differences between single cells.
EigenScore needs five non-empty generations per problem (ten on BCB), so the problem-level
comparison leaves out the problems with fewer: 40 for Qwen3-Coder, none of them
solved, and 143 for gpt-oss, 47 of them solved.

\paragraph{External validity}
We study two open-weight models on competitive-programming problems from LiveCodeBench, on a subset of the same problems in Python, C++ and Java, which
isolates the language, and on more practical Python tasks from BigCodeBench.
Other models, sampling settings or repository-level tasks may behave differently.



\section{Conclusion}
\label{sec:conclusion}
We compared execution-free predictors of code correctness from four families
under benchmark and language shift, for two targets: whether a program is
correct and whether the model's verdict on it can be trusted. Causal
localization, adapted from natural-language question answering, gives the
most robust predictor of the model's verdict under language shift on one of
two models, and its gain lies in the learned subspace, so a simple probe on it
suffices. 
Beyond that, how well a predictor generalizes depends on the target
and the type of shift: no verdict predictor transfers reliably between
benchmarks, while program correctness transfers better. Much of what the predictors
recognize is the problem rather than the candidate, which makes them better suited to deciding which problems to verify than to choosing among candidates. 
Future work should test whether these signals can gate test execution and stronger judges in agentic settings, and whether a few labels from the target setting suffice to adapt them. %, as they do for probes in other domains~\cite{McKenzie2025HighStakes}.



%\newpage

\bibliographystyle{plain}
\bibliography{ood_correctness}
\end{document}
